% =====================================================================
%  Διπλωματική εργασία — κύριο αρχείο
%  Μεταγλώττιση: XeLaTeX + Biber   (latexmk -xelatex main.tex)
%  Overleaf: Menu › Compiler: XeLaTeX
% =====================================================================
\documentclass[a4paper,12pt,oneside,openany]{report}

% Οι σημειώσεις καθοδήγησης (\guide, \todo) κρύβονται με \FinalVersiontrue στο setup.tex.
% =====================================================================
%  Προοίμιο: γλώσσες, γραμματοσειρές, σελίδα, εικόνες, κώδικας, βιβλιογραφία
% =====================================================================

% ---------------- Λειτουργία final -----------------------------------
% \FinalVersiontrue κρύβει τις σημειώσεις καθοδήγησης. Άλλαξέ το στο τέλος.
\newif\ifFinalVersion
\FinalVersionfalse

% ---------------- Γλώσσες και γραμματοσειρές (XeLaTeX) ----------------
\usepackage{fontspec}
\usepackage{polyglossia}
\setmainlanguage{greek}
\setotherlanguage{english}
% Libertinus: πλήρη ελληνικά, μαθηματικά που ταιριάζουν, διαθέσιμη σε TeX Live και Overleaf.
\setmainfont{Libertinus Serif}
\setsansfont{Libertinus Sans}
\setmonofont{DejaVu Sans Mono}[Scale=MatchLowercase]
\newfontfamily\greekfont{Libertinus Serif}
\newfontfamily\greekfontsf{Libertinus Sans}
\newfontfamily\greekfonttt{DejaVu Sans Mono}[Scale=MatchLowercase]

% ---------------- Μαθηματικά -----------------------------------------
\usepackage{amsmath,amssymb}
\usepackage{unicode-math}
\setmathfont{Libertinus Math}
\usepackage{siunitx}
\sisetup{output-decimal-marker={,}, group-separator={.}, detect-all}

% ---------------- Σελίδα ---------------------------------------------
\usepackage[a4paper,top=2.5cm,bottom=2.5cm,left=3cm,right=2.5cm]{geometry}
\usepackage{setspace}
\onehalfspacing
\usepackage{fancyhdr}
\pagestyle{fancy}
\fancyhf{}
\fancyhead[L]{\small\nouppercase{\leftmark}}
\fancyfoot[C]{\thepage}
\renewcommand{\headrulewidth}{0.4pt}
\fancypagestyle{plain}{\fancyhf{}\fancyfoot[C]{\thepage}\renewcommand{\headrulewidth}{0pt}}
\usepackage{titlesec}
\titleformat{\chapter}[display]{\normalfont\huge\bfseries}{\chaptertitlename\ \thechapter}{16pt}{\Huge}
\titlespacing*{\chapter}{0pt}{0pt}{30pt}

% ---------------- Εικόνες, πίνακες, διαγράμματα ----------------------
\usepackage{graphicx}
\graphicspath{{../figures/}{../figures/generated/}}
\usepackage{float}
\usepackage{booktabs,tabularx,multirow,longtable}
\usepackage{adjustbox}
\usepackage[font=small,labelfont=bf]{caption}
\usepackage{subcaption}
\usepackage{tikz}
\usetikzlibrary{arrows.meta,positioning,fit,backgrounds,calc,shapes.geometric}
\usepackage{pgfplots}
\pgfplotsset{compat=1.18}
\usepackage{xcolor}

% ---------------- Κώδικας --------------------------------------------
\usepackage{listings}
\lstdefinestyle{code}{
  basicstyle=\ttfamily\footnotesize,
  keywordstyle=\color{blue!70!black}\bfseries,
  commentstyle=\color{green!40!black}\itshape,
  stringstyle=\color{red!60!black},
  numbers=left, numberstyle=\tiny\color{gray}, numbersep=6pt,
  frame=single, rulecolor=\color{gray!40},
  breaklines=true, showstringspaces=false, tabsize=4,
  captionpos=b, xleftmargin=1.5em, framexleftmargin=1.5em
}
\lstset{style=code}
\renewcommand{\lstlistingname}{Απόσπασμα}
\renewcommand{\lstlistlistingname}{Κατάλογος αποσπασμάτων κώδικα}

% ---------------- Αλγόριθμοι -----------------------------------------
\usepackage{algorithm}
\usepackage{algpseudocode}
\floatname{algorithm}{Αλγόριθμος}
\renewcommand{\listalgorithmname}{Κατάλογος αλγορίθμων}

% ---------------- Βιβλιογραφία ---------------------------------------
\usepackage{csquotes}
\usepackage[backend=biber,style=numeric-comp,sorting=none,maxbibnames=6,giveninits=true]{biblatex}
\addbibresource{../bibliography/references.bib}

% ---------------- Υπερσύνδεσμοι (τελευταίο) --------------------------
\usepackage[hidelinks,unicode]{hyperref}
\usepackage[nameinlink,noabbrev]{cleveref}
\crefname{chapter}{κεφάλαιο}{κεφάλαια}
\Crefname{chapter}{Κεφάλαιο}{Κεφάλαια}
\crefname{section}{ενότητα}{ενότητες}
\Crefname{section}{Ενότητα}{Ενότητες}
\crefname{subsection}{ενότητα}{ενότητες}
\Crefname{subsection}{Ενότητα}{Ενότητες}
\crefname{figure}{σχήμα}{σχήματα}
\Crefname{figure}{Σχήμα}{Σχήματα}
\crefname{table}{πίνακας}{πίνακες}
\Crefname{table}{Πίνακας}{Πίνακες}
\crefname{equation}{εξίσωση}{εξισώσεις}
\Crefname{equation}{Εξίσωση}{Εξισώσεις}
\crefname{lstlisting}{απόσπασμα}{αποσπάσματα}
\Crefname{lstlisting}{Απόσπασμα}{Αποσπάσματα}
\crefname{algorithm}{αλγόριθμος}{αλγόριθμοι}
\Crefname{algorithm}{Αλγόριθμος}{Αλγόριθμοι}

% ---------------- Σημειώσεις καθοδήγησης -----------------------------
% \guide{...}   : τι γράφεται σε αυτή την ενότητα (μπλε πλαίσιο)
% \todo{...}    : κάτι που λείπει (πορτοκαλί πλαίσιο)
% \fillin{...}  : στοιχείο που πρέπει να συμπληρωθεί (κόκκινο κείμενο)
% \code{...}    : όνομα αρχείου / κλάσης του project
\usepackage{tcolorbox}
\ifFinalVersion
  \newcommand{\guide}[1]{}
  \newcommand{\todo}[1]{}
  \newcommand{\fillin}[1]{#1}
\else
  \newtcolorbox{guidebox}{colback=blue!4,colframe=blue!45!black,boxrule=0.4pt,
    left=4pt,right=4pt,top=2pt,bottom=2pt,fontupper=\small\sffamily,
    title={\scriptsize\sffamily ΟΔΗΓΙΑ},fonttitle=\bfseries,coltitle=white}
  \newtcolorbox{todobox}{colback=orange!8,colframe=orange!70!black,boxrule=0.4pt,
    left=4pt,right=4pt,top=2pt,bottom=2pt,fontupper=\small\sffamily,
    title={\scriptsize\sffamily ΕΚΚΡΕΜΕΙ},fonttitle=\bfseries,coltitle=white}
  \newcommand{\guide}[1]{\begin{guidebox}#1\end{guidebox}}
  \newcommand{\todo}[1]{\begin{todobox}#1\end{todobox}}
  \newcommand{\fillin}[1]{\textcolor{red}{[#1]}}
\fi
\newcommand{\code}[1]{\texttt{\detokenize{#1}}}
\newcommand{\en}[1]{\textenglish{#1}}


% ---------------------------------------------------------------------
%  Στοιχεία εργασίας (αλλάζουν μόνο εδώ)
% ---------------------------------------------------------------------
\newcommand{\ThesisTitleEl}{Οπτικοποίηση δεδομένων παρακολούθησης κατασκευών σε μικτή πραγματικότητα: γεωαναφερμένη απόδοση αισθητήρων και τρισδιάστατων μοντέλων σε Meta Quest~3}
\newcommand{\ThesisTitleEn}{Mixed Reality Visualization of Structural Robustness Monitoring Data: Geo-Registered Rendering of Sensors and 3D Models on Meta Quest~3}
\newcommand{\AuthorEl}{Αριστοτέλης Γκούβας}            % έλεγξε την ορθογραφία
\newcommand{\AuthorEn}{Aristotelis Gkouvas}
\newcommand{\StudentId}{\fillin{ΑΜ}}
\newcommand{\SupervisorEl}{\fillin{Ονοματεπώνυμο, βαθμίδα}}
\newcommand{\University}{Πανεπιστήμιο Ιωαννίνων}
\newcommand{\School}{\fillin{Σχολή}}
\newcommand{\Department}{\fillin{Τμήμα}}
\newcommand{\ThesisPlace}{Ιωάννινα}
\newcommand{\ThesisDate}{\fillin{Μήνας} 2027}

\begin{document}

% ---------------- Εισαγωγικές σελίδες (λατινική αρίθμηση) ------------
\pagenumbering{roman}
% Σελίδα τίτλου. Αν το τμήμα δίνει δικό του πρότυπο εξωφύλλου, αντικατέστησε αυτό το αρχείο.
\begin{titlepage}
  \centering
  {\large\scshape \University\par}
  \vspace{0.3cm}
  {\large \School\par}
  {\large \Department\par}
  \vspace{3cm}
  {\Large\bfseries Διπλωματική Εργασία\par}
  \vspace{1.2cm}
  {\LARGE\bfseries \ThesisTitleEl\par}
  \vspace{0.35cm}
  {\Large \ThesisSubtitleEl\par}
  \vspace{0.8cm}
  {\large\itshape\selectlanguage{english} \ThesisTitleEn\\ \ThesisSubtitleEn\par}
  \vspace{2.5cm}
  {\Large \AuthorEl\par}
  \vspace{0.2cm}
  {\normalsize ΑΜ: \StudentId\par}
  \vfill
  {\large Επιβλέπων: \SupervisorEl\par}
  \vspace{1.5cm}
  {\large \ThesisPlace, \ThesisDate\par}
\end{titlepage}

\thispagestyle{empty}
\vspace*{\fill}
\begin{center}
  \textbf{Τριμελής εξεταστική επιτροπή}\\[1em]
  \begin{tabular}{l}
    \SupervisorEl\ (επιβλέπων) \\
    \fillin{Μέλος επιτροπής 2} \\
    \fillin{Μέλος επιτροπής 3} \\
  \end{tabular}
\end{center}
\vspace{2cm}
\noindent\small
\guide{Πολλά τμήματα απαιτούν εδώ και τη δήλωση μη λογοκλοπής / πνευματικών δικαιωμάτων. Πάρε το ακριβές κείμενο από τη γραμματεία ή το πρότυπο του τμήματος.}
\vspace*{\fill}
\clearpage

\chapter*{Ευχαριστίες}
\addcontentsline{toc}{chapter}{Ευχαριστίες}

\guide{Γράφεται στο τέλος. Επιβλέπων, όσοι έδωσαν πρόσβαση σε εργοτάξιο / δεδομένα Terraweb, όσοι συμμετείχαν στη μελέτη χρηστών, οικογένεια. Μισή σελίδα αρκεί.}

\chapter*{Περίληψη}
\addcontentsline{toc}{chapter}{Περίληψη}

\guide{Γράφεται τελευταία, όταν υπάρχουν αριθμοί από το \cref{ch:evaluation}. 250--350 λέξεις, μία παράγραφος ή τέσσερις μικρές:
(1) πρόβλημα, (2) τι κατασκευάστηκε, (3) πώς αξιολογήθηκε, (4) τα κύρια αποτελέσματα \emph{με αριθμούς} (π.χ. «οριζόντιο σφάλμα \fillin{x}\,cm σε απόσταση \fillin{y}\,m από τα σημεία ελέγχου», «σταθεροί 72\,fps έως \fillin{z} χιλιάδες τρίγωνα», «SUS \fillin{w}»).}

Οι μετρήσεις των αισθητήρων παρακολούθησης κατασκευών απεικονίζονται συνήθως σε οθόνες, μακριά από το σημείο όπου λαμβάνονται, και ο μηχανικός στο πεδίο πρέπει να αντιστοιχίσει νοερά κάθε σημείο του χάρτη με τη θέση του στην κατασκευή. Η εργασία παρουσιάζει ένα end-to-end σύστημα που μεταφέρει τις μετρήσεις από μια πλατφόρμα παρακολούθησης IoT σε ένα αυτόνομο headset μικτής πραγματικότητας (Meta Quest~3) και τις εμφανίζει πάνω στους ίδιους τους αισθητήρες, στον πραγματικό χώρο.
\fillin{Συνέχεια: ευθυγράμμιση, απεικονίσεις, αξιολόγηση, αποτελέσματα.}

\vspace{1em}
\noindent\textbf{Λέξεις-κλειδιά:} μικτή πραγματικότητα, παρακολούθηση κατασκευών, διαδίκτυο των πραγμάτων, γεωαναφορά, χωρική ευθυγράμμιση, Meta Quest~3, οπτικοποίηση δεδομένων

\begin{otherlanguage}{english}
\chapter*{Abstract}
\addcontentsline{toc}{chapter}{Abstract}

\guide{The English abstract is a translation of the Greek one, not a different text.}

Measurements from structural monitoring sensors are usually viewed on screens, away from where they are taken, so an engineer on site has to map every point on the plan to its place on the structure. This thesis presents an end-to-end system that carries measurements from an IoT monitoring platform to a standalone mixed reality headset (Meta Quest~3) and shows them on the sensors themselves, in physical space.
\fillin{Continued.}

\vspace{1em}
\noindent\textbf{Keywords:} mixed reality, structural health monitoring, Internet of Things, georeferencing, spatial registration, Meta Quest~3, data visualization
\end{otherlanguage}


\tableofcontents
\listoffigures
\listoftables
\chapter*{Συντομογραφίες}
\addcontentsline{toc}{chapter}{Συντομογραφίες}

\begin{tabularx}{\textwidth}{@{}lX@{}}
\toprule
\textbf{Συντομογραφία} & \textbf{Ανάπτυξη} \\
\midrule
API   & \en{Application Programming Interface} \\
AR    & Επαυξημένη πραγματικότητα (\en{Augmented Reality}) \\
BIM   & \en{Building Information Modeling} \\
BiRP  & \en{Built-in Render Pipeline} (Unity) \\
CRS   & Σύστημα αναφοράς συντεταγμένων (\en{Coordinate Reference System}) \\
CSV   & \en{Comma-Separated Values} \\
DXF   & \en{Drawing Exchange Format} \\
ΕΓΣΑ87 & Ελληνικό Γεωδαιτικό Σύστημα Αναφοράς 1987 \\
FPS   & Καρέ ανά δευτερόλεπτο (\en{Frames Per Second}) \\
HTTP  & \en{Hypertext Transfer Protocol} \\
IDW   & \en{Inverse Distance Weighting} \\
IFC   & \en{Industry Foundation Classes} \\
IoT   & Διαδίκτυο των πραγμάτων (\en{Internet of Things}) \\
LAN   & Τοπικό δίκτυο (\en{Local Area Network}) \\
MR    & Μικτή πραγματικότητα (\en{Mixed Reality}) \\
NASA-TLX & \en{NASA Task Load Index} \\
OBJ   & \en{Wavefront OBJ} \\
QEM   & \en{Quadric Error Metrics} \\
REST  & \en{Representational State Transfer} \\
RMS(E) & Ρίζα μέσου τετραγωνικού (σφάλματος) (\en{Root Mean Square (Error)}) \\
SHM   & Παρακολούθηση δομικής ακεραιότητας (\en{Structural Health Monitoring}) \\
SLAM  & \en{Simultaneous Localization and Mapping} \\
SSQ   & \en{Simulator Sickness Questionnaire} \\
SUS   & \en{System Usability Scale} \\
VR    & Εικονική πραγματικότητα (\en{Virtual Reality}) \\
XR    & Εκτεταμένη πραγματικότητα (\en{Extended Reality}) \\
\bottomrule
\end{tabularx}


% ---------------- Κυρίως κείμενο --------------------------------------
\cleardoublepage
\pagenumbering{arabic}
\chapter{Εισαγωγή}
\label{ch:intro}

\section{Το πρόβλημα}
\label{sec:intro-problem}

Η παρακολούθηση της δομικής ακεραιότητας (\en{Structural Health Monitoring}, SHM) γεφυρών, φραγμάτων, σηράγγων, πρανών και κτιρίων βασίζεται σε δίκτυα αισθητήρων που καταγράφουν συνεχώς μετακινήσεις, κλίσεις, ανοίγματα ρωγμών, πιέσεις και άλλα μεγέθη \cite{farrar2007shm}. Σε ένα τυπικό έργο, δεκάδες ή εκατοντάδες σημεία μέτρησης (κλισιόμετρα, ρωγμόμετρα, πιεζόμετρα, τοπογραφικοί στόχοι) τροφοδοτούν μια πλατφόρμα παρακολούθησης στο διαδίκτυο. Εκεί οι μετρήσεις αποθηκεύονται ως χρονοσειρές, συγκρίνονται με όρια συναγερμού και παρουσιάζονται ως πίνακες, διαγράμματα και κατόψεις.

Η πληροφορία όμως αφορά συγκεκριμένα σημεία του φυσικού χώρου, ενώ η παρουσίασή της γίνεται μακριά από αυτά, σε μια οθόνη. Ο μηχανικός που επισκέπτεται το έργο πρέπει να κάνει νοερά την αντιστοίχιση: ο αισθητήρας «Κ7» της λίστας, που η πλατφόρμα δείχνει σε κατάσταση συναγερμού, βρίσκεται σε κάποιο σημείο της κατόψεως, και το σημείο αυτό σε κάποια θέση του πραγματικού τοίχου, βάθρου ή πρανούς μπροστά του. Η αντιστοίχιση αυτή κοστίζει χρόνο, γίνεται συχνά λάθος σε έργα με πολλούς παρόμοιους αισθητήρες, και αφήνει αναξιοποίητη τη χωρική πληροφορία που έχει ο ίδιος ο χώρος: ποιοι αισθητήρες κινούνται μαζί, αν μια ζώνη μετακίνησης συμπίπτει με μια ορατή ρωγμή, πού βρίσκεται ο επόμενος αισθητήρας προς επιθεώρηση.

\review{Η παράγραφος περιγράφει τη ροή εργασίας «από γενική εμπειρία». Πρόσθεσε ένα \textbf{συγκεκριμένο παράδειγμα από δικό σου έργο ή επίσκεψη} (π.χ. το έργο όπου δοκιμάστηκε η εφαρμογή, πόσοι αισθητήρες, πόσος χρόνος χρειάζεται σήμερα για μια επιθεώρηση). Αν δεν υπάρχει μέτρηση, μην βάλεις αριθμούς.}

\section{Κίνητρο}
\label{sec:intro-motivation}

Η επαυξημένη και η μικτή πραγματικότητα (AR/MR) προσφέρουν μια άμεση απάντηση στο πρόβλημα: να εμφανίζεται η πληροφορία \emph{πάνω} στο σημείο που αφορά, στο οπτικό πεδίο του χρήστη \cite{azuma1997survey,bressa2022situated}. Η ιδέα δεν είναι νέα, αλλά μέχρι πρόσφατα προσέκρουε σε πρακτικά εμπόδια: ακριβές συσκευές οπτικής διαφάνειας με μικρό οπτικό πεδίο, ανάγκη για υπολογιστή ή καλώδια, και κυρίως η δυσκολία να τοποθετηθεί με ακρίβεια εκατοστών ένα γεωαναφερμένο σημείο στον χώρο μιας συσκευής που δεν γνωρίζει το γεωδαιτικό σύστημα αναφοράς.

Τα αυτόνομα headsets της τελευταίας γενιάς, όπως το Meta Quest~3, αλλάζουν την εικόνα. Συνδυάζουν έγχρωμο passthrough υψηλής ανάλυσης, εντοπισμό θέσης χωρίς εξωτερικό εξοπλισμό, αισθητήρα βάθους και spatial anchors που επιβιώνουν μεταξύ συνεδριών, με κόστος συγκρίσιμο με ένα tablet. Αυτό κάνει ρεαλιστική μια εφαρμογή πεδίου που ο μηχανικός μπορεί να φορέσει στο εργοτάξιο. Ανοιχτά παραμένουν τρία ερωτήματα: πόσο ακριβής μπορεί να είναι η ευθυγράμμιση με απλά μέσα, πόσο σταθερή μένει, και αν η εμπειρία βοηθά πραγματικά τον χρήστη.

\section{Ερευνητικά ερωτήματα}
\label{sec:rq}

Η εργασία εξετάζει τα εξής ερωτήματα:

\begin{description}
  \item[ΕΕ1 (ακρίβεια).] Με ποια ακρίβεια μπορεί ένα αυτόνομο headset μικτής πραγματικότητας να τοποθετήσει γεωαναφερμένα σημεία στη θέση τους στον πραγματικό χώρο, με λίγα σημεία ελέγχου και χωρίς εξωτερικό εξοπλισμό εντοπισμού;
  \item[ΕΕ2 (σταθερότητα).] Πόσο σταθερή μένει η ευθυγράμμιση στον χρόνο και μεταξύ συνεδριών μέσω spatial anchors, και πώς επηρεάζεται από τον φωτισμό και τις αλλαγές στο εργοτάξιο;
  \item[ΕΕ3 (απόδοση).] Ποια είναι τα όρια φόρτου (τρίγωνα μοντέλου, πλήθος αισθητήρων) για σταθερό ρυθμό καρέ σε αυτόνομη συσκευή, και πόσο τα μετακινούν οι τεχνικές βελτιστοποίησης;
  \item[ΕΕ4 (χρηστικότητα).] Βελτιώνει η επιτόπια απεικόνιση την ταχύτητα και την ακρίβεια μιας επιθεώρησης σε σχέση με τη συμβατική ροή εργασίας;
\end{description}

\section{Συνεισφορές}
\label{sec:contributions}

\begin{enumerate}
  \item Σχεδίαση και υλοποίηση ενός end-to-end συστήματος: πλατφόρμα παρακολούθησης (Terraweb) → εργαλεία PC σε Python → ενσωματωμένος εξυπηρετητής HTTP στη συσκευή → απεικόνιση σε μικτή πραγματικότητα (\cref{ch:design,ch:implementation}).
  \item Μέθοδος ευθυγράμμισης γεωαναφερμένου έργου με τον χώρο της συσκευής, με κλειστή λύση ελαχίστων τετραγώνων και πρόβλεψη σφάλματος σε πραγματικό χρόνο που καθοδηγεί τον χρήστη στο επόμενο σημείο ελέγχου (\cref{sec:impl-alignment}).
  \item Τεχνικές απεικόνισης δεδομένων αισθητήρων στον χώρο: κορδέλες ιστορικού, χρονομηχανή, διανύσματα μετακίνησης, χάρτης θερμότητας με παρεμβολή πάνω σε τριγωνοποίηση Delaunay, ράβδοι εγγύτητας στα όρια, ζωντανοί συναγερμοί, διαδρομή επιθεώρησης (\cref{sec:impl-vis}).
  \item Αλυσίδα προετοιμασίας και φόρτωσης βαριών μοντέλων (φωτογραμμετρία, DXF, IFC) σε συσκευή περιορισμένων πόρων (\cref{sec:impl-models}).
  \item Πειραματική αξιολόγηση ακρίβειας, σταθερότητας, απόδοσης και χρηστικότητας στο πεδίο (\cref{ch:evaluation}).
\end{enumerate}

\review[!]{Τα ΕΕ και οι συνεισφορές πρέπει στο τέλος να \textbf{ταιριάζουν με όσα δείχνει πράγματι το κεφ. 6}. Αν ένα πείραμα δεν γίνει (π.χ. η μελέτη χρηστών), αφαίρεσε ή μαλάκωσε το αντίστοιχο ΕΕ και τη συνεισφορά 5. Συζήτησέ τα με τον επιβλέποντα πριν ξεκινήσουν οι μετρήσεις.}

\section{Δομή της εργασίας}
\label{sec:outline}

Το \cref{ch:background} παρουσιάζει το θεωρητικό υπόβαθρο: μικτή πραγματικότητα, συστήματα συντεταγμένων, χωρική ευθυγράμμιση, παρεμβολή και απόδοση γραφικών σε κινητές συσκευές. Το \cref{ch:related} εξετάζει σχετικές εργασίες. Το \cref{ch:design} περιγράφει τις απαιτήσεις και την αρχιτεκτονική του συστήματος, και το \cref{ch:implementation} την υλοποίηση των υποσυστημάτων του. Το \cref{ch:evaluation} παρουσιάζει την πειραματική αξιολόγηση, το \cref{ch:discussion} συζητά τα αποτελέσματα και τους περιορισμούς, και το \cref{ch:conclusions} κλείνει με συμπεράσματα και μελλοντική εργασία.

\chapter{Θεωρητικό υπόβαθρο}
\label{ch:background}

Το κεφάλαιο συγκεντρώνει τις έννοιες στις οποίες στηρίζονται τα επόμενα: τι είναι η μικτή πραγματικότητα και τι προσφέρει η συσκευή που χρησιμοποιήθηκε, πώς περιγράφονται οι θέσεις των αισθητήρων, πώς ευθυγραμμίζονται δύο συστήματα συντεταγμένων, πώς εκτιμάται ένα μέγεθος ανάμεσα σε μετρήσεις, ποιοι περιορισμοί ισχύουν στην απόδοση γραφικών σε αυτόνομες συσκευές, και πώς οργανώνεται η παρακολούθηση κατασκευών.

\section{Μικτή πραγματικότητα}
\label{sec:bg-mr}

Οι \en{Milgram} και \en{Kishino} περιέγραψαν το \emph{συνεχές πραγματικότητας--εικονικότητας}, που εκτείνεται από το εντελώς πραγματικό περιβάλλον ως το εντελώς εικονικό \cite{milgram1994taxonomy}. Ενδιάμεσα βρίσκεται η \emph{μικτή πραγματικότητα} (MR), όπου πραγματικά και εικονικά στοιχεία συνυπάρχουν. Η επαυξημένη πραγματικότητα (AR) είναι η περιοχή όπου κυριαρχεί το πραγματικό περιβάλλον και τα εικονικά στοιχεία προστίθενται πάνω του.

Ο \en{Azuma} όρισε ένα σύστημα AR με τρία χαρακτηριστικά: συνδυάζει πραγματικό και εικονικό, είναι διαδραστικό σε πραγματικό χρόνο, και τα εικονικά στοιχεία είναι \emph{ευθυγραμμισμένα} (\en{registered}) στον τρισδιάστατο χώρο \cite{azuma1997survey}. Το τρίτο χαρακτηριστικό είναι και το δυσκολότερο: ένα εικονικό σημείο που «κολυμπά» μερικές δεκάδες εκατοστά μακριά από τον αισθητήρα που αντιπροσωπεύει μπορεί να παραπλανήσει περισσότερο από ό,τι βοηθά. Για τον λόγο αυτό η ευθυγράμμιση είναι κεντρικό θέμα της παρούσας εργασίας. Αναλυτική επισκόπηση του πεδίου δίνουν οι \en{Billinghurst} κ.ά. \cite{billinghurst2015survey}.

Οι συσκευές κεφαλής χωρίζονται σε δύο κατηγορίες. Στις συσκευές \emph{οπτικής διαφάνειας} (\en{optical see-through}, π.χ. HoloLens) ο χρήστης βλέπει τον πραγματικό κόσμο μέσα από ημιδιαφανή οπτικά στοιχεία, πάνω στα οποία προβάλλονται τα εικονικά. Στις συσκευές \emph{video passthrough} (π.χ. Quest~3) κάμερες καταγράφουν τον περιβάλλοντα χώρο και η εικόνα του αναπαράγεται στις οθόνες μαζί με τα εικονικά στοιχεία. Οι δεύτερες έχουν μεγαλύτερο οπτικό πεδίο και πλήρη έλεγχο της σύνθεσης (π.χ. μπορούν να σκοτεινιάσουν τον πραγματικό χώρο), με κόστος την ανάλυση και την καθυστέρηση της κάμερας.

\subsection{Η συσκευή Meta Quest 3}
\label{sec:bg-quest}

Το Meta Quest~3 είναι αυτόνομο headset με επεξεργαστή Qualcomm Snapdragon XR2 Gen~2 και λειτουργικό βασισμένο στο Android (Horizon OS). Για την παρούσα εργασία έχουν σημασία τα εξής χαρακτηριστικά:
\begin{itemize}
  \item \textbf{Εντοπισμός θέσης inside-out}: η συσκευή εκτιμά τη θέση και τον προσανατολισμό της με έξι βαθμούς ελευθερίας από τις κάμερες και τους αδρανειακούς αισθητήρες της (SLAM), χωρίς εξωτερικούς σταθμούς. Η κατακόρυφος είναι γνωστή με ακρίβεια από το επιταχυνσιόμετρο.
  \item \textbf{Έγχρωμο passthrough} από δύο κάμερες RGB, που επιτρέπει εφαρμογές μικτής πραγματικότητας.
  \item \textbf{Αισθητήρας βάθους}, μέσω του οποίου η εφαρμογή γνωρίζει την απόσταση των πραγματικών επιφανειών και μπορεί να «κρύψει» εικονικά αντικείμενα πίσω τους (απόκρυψη, \en{occlusion}).
  \item \textbf{Χειριστήρια} με εντοπισμό έξι βαθμών ελευθερίας, που χρησιμοποιούνται εδώ και ως όργανο σκόπευσης: η άκρη τους ακουμπά τα σημεία ελέγχου.
  \item \textbf{Ρυθμός ανανέωσης} 72, 90 ή 120\,Hz, και \textbf{foveated rendering} σταθερού τύπου.
\end{itemize}

\review{Επιβεβαίωσε τα τεχνικά χαρακτηριστικά του Quest~3 από \textbf{επίσημη πηγή} (σελίδα προϊόντος ή τεκμηρίωση της Meta) και πρόσθεσε αναφορά (\code{@online} στο \code{references.bib} με ημερομηνία πρόσβασης). Αν ο επιβλέπων θέλει πίνακα χαρακτηριστικών (ανάλυση, οπτικό πεδίο, μνήμη), πρόσθεσέ τον εδώ.}

\subsection{Spatial anchors}
\label{sec:bg-anchors}

Ο εντοπισμός θέσης της συσκευής ορίζει ένα σύστημα συντεταγμένων που ισχύει μόνο για την τρέχουσα συνεδρία: η αρχή του μπορεί να αλλάξει σε κάθε εκκίνηση ή επανακέντρωση. Για να «κολλήσει» ένα εικονικό στοιχείο σε ένα σημείο του πραγματικού χώρου με διάρκεια, οι πλατφόρμες προσφέρουν \emph{spatial anchors}: σημεία που η συσκευή συνδέει με τον χάρτη οπτικών χαρακτηριστικών του περιβάλλοντος. Ένα anchor αποθηκεύεται με ένα αναγνωριστικό (UUID) και, σε επόμενη συνεδρία, η συσκευή προσπαθεί να το \emph{επανεντοπίσει} (\en{relocalize}) αναγνωρίζοντας τον χώρο. Η επιτυχία και η ακρίβεια του επανεντοπισμού εξαρτώνται από τον φωτισμό, την υφή των επιφανειών και τις αλλαγές του χώρου από την προηγούμενη επίσκεψη. Επιπλέον, η ακρίβεια ενός anchor υποβαθμίζεται με την απόσταση από αυτό, κάτι που σε ένα μεγάλο εργοτάξιο οδηγεί στη χρήση περισσότερων του ενός anchors. Η συμπεριφορά αυτή μετριέται πειραματικά στην \cref{sec:exp-anchors}.

\section{Συστήματα συντεταγμένων}
\label{sec:bg-crs}

Οι θέσεις των αισθητήρων ενός έργου δίνονται από τοπογραφική αποτύπωση σε ένα \emph{προβολικό} σύστημα συντεταγμένων: ένα επίπεδο, μετρικό σύστημα με άξονες Ανατολή ($E$) και Βορράς ($N$), και ένα υψόμετρο ($H$). Στην Ελλάδα το συνηθέστερο είναι το ΕΓΣΑ87, ενώ σε άλλα έργα χρησιμοποιούνται ζώνες UTM ή τοπικοί κάναβοι του έργου. Το σύστημα που περιγράφεται εδώ δεν δένεται σε συγκεκριμένο σύστημα: δέχεται οποιονδήποτε μετρικό κάναβο $E/N/H$, αρκεί αισθητήρες, σημεία ελέγχου και μοντέλα να εκφράζονται στον ίδιο.

Οι μηχανές απόδοσης γραφικών χρησιμοποιούν δικό τους σύστημα. Η Unity έχει αριστερόστροφο σύστημα με τον άξονα $y$ προς τα πάνω, ενώ ο κάναβος του έργου είναι δεξιόστροφος με το $H$ προς τα πάνω. Η αντιστοίχιση από τον κάναβο $(E, N, H)$ στο τοπικό σύστημα $(x, y, z)$ είναι
\begin{equation}
  \begin{pmatrix} x \\ y \\ z \end{pmatrix} =
  \begin{pmatrix} E - E_0 \\ H - H_0 \\ N - N_0 \end{pmatrix},
  \label{eq:enh-to-local}
\end{equation}
όπου $(E_0, N_0, H_0)$ η αρχή του έργου. Η αφαίρεση της αρχής είναι απαραίτητη: οι συντεταγμένες ΕΓΣΑ87 έχουν τάξη μεγέθους $10^5$--$10^6$\,m, και σε αριθμούς κινητής υποδιαστολής απλής ακρίβειας (24 bit mantissa) αυτό αφήνει ακρίβεια λίγων εκατοστών. Με την αρχή μέσα στο έργο, οι τιμές μένουν μικρές και η ακρίβεια φτάνει το χιλιοστό. Η αρχή κρατιέται σε διπλή ακρίβεια.

\section{Χωρική ευθυγράμμιση}
\label{sec:bg-registration}

\emph{Ευθυγράμμιση} (\en{registration}) είναι η εκτίμηση του μετασχηματισμού που φέρνει ένα σύνολο σημείων από ένα σύστημα συντεταγμένων σε ένα άλλο. Όταν είναι γνωστές αντιστοιχίες σημείων $(\mathbf{p}_i, \mathbf{m}_i)$, $i = 1, \dots, n$, το πρόβλημα είναι γνωστό ως πρόβλημα του \emph{απόλυτου προσανατολισμού}: ζητούνται στροφή $R$, μετάθεση $\mathbf{t}$ και, προαιρετικά, κλίμακα $s$ που ελαχιστοποιούν
\begin{equation}
  E(R, \mathbf{t}, s) = \sum_{i=1}^{n} \lVert s R \mathbf{p}_i + \mathbf{t} - \mathbf{m}_i \rVert^2 .
  \label{eq:absolute-orientation}
\end{equation}
Το πρόβλημα έχει κλειστές λύσεις. Ο \en{Horn} τη διατύπωσε με μοναδιαία τετραδόνια \cite{horn1987closed}. Οι \en{Arun} κ.ά. και ο \en{Kabsch} μέσω της ανάλυσης ιδιαζουσών τιμών (SVD) του πίνακα συνδιακύμανσης των κεντραρισμένων σημείων \cite{arun1987leastsquares,kabsch1976solution}. Ο \en{Umeyama} έδωσε τη γενική λύση με κλίμακα και διόρθωσε την περίπτωση όπου η SVD επιστρέφει κατοπτρισμό αντί για στροφή \cite{umeyama1991leastsquares}. Σε όλες τις περιπτώσεις η μετάθεση προκύπτει από τα κέντρα βάρους: $\mathbf{t} = \bar{\mathbf{m}} - sR\bar{\mathbf{p}}$.

Η ποιότητα μιας ευθυγράμμισης αποτιμάται με τα \emph{υπόλοιπα} $\mathbf{r}_i = sR\mathbf{p}_i + \mathbf{t} - \mathbf{m}_i$ και τη ρίζα του μέσου τετραγώνου τους (RMS). Το RMS στα ίδια τα σημεία ευθυγράμμισης είναι όμως αισιόδοξο, γιατί η λύση έχει προσαρμοστεί σε αυτά. Η πραγματική ακρίβεια μετριέται σε \emph{ανεξάρτητα σημεία ελέγχου} (\en{check points}), που δεν συμμετέχουν στη λύση, όπως γίνεται στην τοπογραφία. Αυτή η διάκριση ακολουθείται στο πείραμα της \cref{sec:exp-alignment}.

Όταν δεν υπάρχουν γνωστές αντιστοιχίες αλλά δύο νέφη σημείων ή επιφάνειες, χρησιμοποιείται ο επαναληπτικός αλγόριθμος ICP \cite{besl1992icp}, που εναλλάσσει την εύρεση πλησιέστερων σημείων με την κλειστή λύση της \cref{eq:absolute-orientation}. Η εργασία δεν τον χρησιμοποιεί, αλλά τον αναφέρει ως μελλοντική επέκταση (\cref{sec:future}).

\section{Παρεμβολή χωρικών δεδομένων}
\label{sec:bg-interpolation}

Οι αισθητήρες μετρούν σε διακριτά σημεία. Για να απεικονιστεί ένα μέγεθος, π.χ. η μετακίνηση, πάνω σε μια συνεχή επιφάνεια, πρέπει να εκτιμηθεί η τιμή του ανάμεσα στα σημεία μέτρησης.

Η μέθοδος \emph{αντίστροφης στάθμισης απόστασης} (IDW) του \en{Shepard} εκτιμά την τιμή σε ένα σημείο $\mathbf{x}$ ως σταθμισμένο μέσο όρο όλων των μετρήσεων $z_i$ \cite{shepard1968interpolation}:
\begin{equation}
  \hat{z}(\mathbf{x}) = \frac{\sum_i w_i(\mathbf{x})\, z_i}{\sum_i w_i(\mathbf{x})}, \qquad
  w_i(\mathbf{x}) = \lVert \mathbf{x} - \mathbf{x}_i \rVert^{-p},
  \label{eq:idw}
\end{equation}
με εκθέτη $p$ συνήθως ίσο με 2. Είναι απλή και ορίζεται παντού, αλλά δημιουργεί χαρακτηριστικά «μάτια» γύρω από κάθε μέτρηση, γιατί η επιφάνεια επιπεδώνεται κοντά σε κάθε σημείο, και επηρεάζεται από μακρινά σημεία.

Η \emph{τριγωνοποίηση Delaunay} ενός συνόλου σημείων στο επίπεδο είναι η τριγωνοποίηση όπου κανένα σημείο δεν βρίσκεται μέσα στον περιγεγραμμένο κύκλο κάποιου τριγώνου. Μεγιστοποιεί την ελάχιστη γωνία και αποφεύγει τα στενόμακρα τρίγωνα \cite{deberg2008computational}. Πάνω σε αυτήν, η \emph{γραμμική παρεμβολή} δίνει σε κάθε σημείο τον βαρυκεντρικό συνδυασμό των τιμών των τριών κορυφών του τριγώνου που το περιέχει. Το αποτέλεσμα είναι συνεχές, περνά ακριβώς από τις μετρήσεις και χρησιμοποιεί μόνο τους γειτονικούς αισθητήρες. Μειονέκτημα είναι ότι δεν ορίζεται έξω από το κυρτό περίβλημα των σημείων. Πιο σύνθετες γεωστατιστικές μέθοδοι, όπως το \en{kriging}, μοντελοποιούν τη χωρική συσχέτιση αλλά απαιτούν περισσότερα δεδομένα για την εκτίμηση του βαριογράμματος.

\section{Απόδοση γραφικών σε αυτόνομες συσκευές}
\label{sec:bg-rendering}

Μια εφαρμογή εικονικής ή μικτής πραγματικότητας πρέπει να διατηρεί σταθερό ρυθμό καρέ ίσο με τον ρυθμό ανανέωσης της οθόνης. Αλλιώς η εικόνα τρεμοπαίζει και ο χρήστης μπορεί να νιώσει δυσφορία. Στα 72\,Hz κάθε καρέ έχει προϋπολογισμό \SI{13.9}{ms}, και στα 90\,Hz \SI{11.1}{ms}, για δύο εικόνες (μία ανά μάτι).

Σε μια κινητή GPU το κόστος ενός καρέ καθορίζεται κυρίως από τρεις παράγοντες. Ο πρώτος είναι το \emph{πλήθος των draw calls}, δηλαδή των ξεχωριστών εντολών σχεδίασης που στέλνει η CPU. Μειώνεται με την ομαδοποίηση (\en{batching}) αντικειμένων που μοιράζονται υλικό, και με τη σχεδίαση και των δύο ματιών σε ένα πέρασμα (\en{single-pass instanced}). Ο δεύτερος είναι το \emph{πλήθος των τριγώνων και των κορυφών}. Ο τρίτος είναι το \emph{κόστος ανά pixel}, που αυξάνεται απότομα με την ημιδιαφάνεια, γιατί κάθε ημιδιαφανής επιφάνεια ζωγραφίζεται πάνω από τις προηγούμενες (\en{overdraw}). Το \emph{foveated rendering} μειώνει την ανάλυση στην περιφέρεια του οπτικού πεδίου, όπου το ανθρώπινο μάτι διακρίνει λιγότερη λεπτομέρεια, και ελευθερώνει χρόνο GPU \cite{patney2016foveated}.

Τα μοντέλα φωτογραμμετρίας έχουν συχνά εκατομμύρια τρίγωνα και υφές υψηλής ανάλυσης, πολύ πάνω από όσα αντέχει μια αυτόνομη συσκευή. Η \emph{απλοποίηση πλέγματος} μειώνει τα τρίγωνα με τη μικρότερη δυνατή γεωμετρική αλλοίωση. Η μέθοδος των \en{Garland} και \en{Heckbert} συμπτύσσει διαδοχικά ακμές, επιλέγοντας κάθε φορά αυτήν με το μικρότερο τετραγωνικό σφάλμα (QEM) ως προς τα επίπεδα των γειτονικών τριγώνων \cite{garland1997qem}. Η απόκλιση του απλοποιημένου πλέγματος από το αρχικό μετριέται συνήθως με την \emph{απόσταση Hausdorff}, όπως στο εργαλείο \en{Metro} \cite{cignoni1998metro}. Και τα δύο είναι διαθέσιμα στο MeshLab \cite{cignoni2008meshlab}.

\section{Παρακολούθηση κατασκευών και πλατφόρμες IoT}
\label{sec:bg-shm}

Η παρακολούθηση δομικής ακεραιότητας είναι η διαδικασία εντοπισμού μεταβολών σε μια κατασκευή μέσω συστηματικών μετρήσεων στον χρόνο \cite{farrar2007shm}. Σε γεωτεχνικά έργα και υποδομές χρησιμοποιούνται, μεταξύ άλλων:
\begin{itemize}
  \item \textbf{τοπογραφικοί στόχοι} (πρίσματα), που μετρώνται με γεωδαιτικό σταθμό και δίνουν θέση $(x, y, z)$ και μετακινήσεις $(dx, dy, dz)$ ως προς μια μηδενική μέτρηση·
  \item \textbf{κλισιόμετρα}, για τη μεταβολή κλίσης·
  \item \textbf{ρωγμόμετρα}, για το άνοιγμα ρωγμών και αρμών·
  \item \textbf{πιεζόμετρα}, για τη στάθμη ή την πίεση του νερού.
\end{itemize}
Σε κάθε αισθητήρα ορίζονται \emph{όρια} σε βαθμίδες (π.χ. αναφορά, προειδοποίηση, συναγερμός). Όταν μια μέτρηση ξεπερνά ένα όριο, η πλατφόρμα δημιουργεί συμβάν συναγερμού.

Οι μετρήσεις συλλέγονται σε μια πλατφόρμα παρακολούθησης, που στην παρούσα εργασία είναι το Terraweb. Οργανώνει τα δεδομένα σε έργα, περιοχές, αισθητήρες, χρονοσειρές μετρήσεων, όρια και ενεργά συμβάντα, και τα εκθέτει μέσω διεπαφής REST \cite{fielding2000rest} με αυθεντικοποίηση JWT. Η πρόσβαση της συσκευής σε αυτή τη διεπαφή περιγράφεται στην \cref{sec:comm-terraweb}.

\review{Επιβεβαίωσε την περιγραφή του Terraweb (οργάνωση σε έργα/περιοχές, όρια σε βαθμίδες, ορολογία των βαθμίδων). Αν η πλατφόρμα έχει δημόσια τεκμηρίωση ή είναι προϊόν εταιρείας, δες με τον επιβλέποντα \textbf{πώς επιτρέπεται να αναφερθεί} (όνομα, εταιρεία, εμπιστευτικότητα).}

\chapter{Σχετικές εργασίες}
\label{ch:related}

\review[!]{\textbf{Το κεφάλαιο χρειάζεται τη μεγαλύτερη δουλειά από σένα.} Οι εργασίες βρέθηκαν με αναζήτηση και οι περιγραφές τους βασίζονται σε περιλήψεις, όχι σε ανάγνωση του πλήρους κειμένου. Πριν κρατήσεις οποιαδήποτε πρόταση:
(1) \textbf{διάβασε κάθε εργασία} (τουλάχιστον περίληψη, μέθοδο, αποτελέσματα) και διόρθωσε όσα γράφω γι' αυτήν·
(2) επιβεβαίωσε τα στοιχεία στο \code{references.bib} (συγγραφείς, τόμος, σελίδες, DOI)·
(3) \textbf{πρόσθεσε περισσότερες εργασίες}: μια διπλωματική θέλει συνήθως 25--40 αναφορές συνολικά. Λέξεις-κλειδιά στον \code{Οδηγός_Διπλωματικής.md} §5.}

Η εργασία βρίσκεται στην τομή τεσσάρων περιοχών: της μικτής πραγματικότητας για επιθεώρηση και παρακολούθηση κατασκευών, του BIM στο εργοτάξιο, της γεωαναφοράς σε επαυξημένη πραγματικότητα, και της οπτικοποίησης δεδομένων στη θέση τους (\en{situated visualization}). Το κεφάλαιο εξετάζει καθεμιά και κλείνει με μια σύγκριση.

\section{Μικτή πραγματικότητα για επιθεώρηση και παρακολούθηση κατασκευών}
\label{sec:rw-shm}

Η ιδέα να συνδεθεί η οπτική επιθεώρηση με ψηφιακά εργαλεία μέσα από συσκευές AR έχει διατυπωθεί ρητά στη βιβλιογραφία της παρακολούθησης κατασκευών. Οι \en{Mascareñas} κ.ά. προτείνουν ένα πλαίσιο λογισμικού για εφαρμογές επιθεώρησης υποδομών με AR. Ξεκινούν από το ότι η οπτική επιθεώρηση παραμένει η κύρια μέθοδος, και δείχνουν πώς μια συσκευή κεφαλής μπορεί να καταγράφει τρισδιάστατες μετρήσεις και να τις συνδέει με την κατασκευή \cite{mascarenas2021ar}.

Πιο κοντά στην παρούσα εργασία είναι συστήματα που φέρνουν \emph{δεδομένα αισθητήρων} στο πεδίο. Οι \en{Fawad} κ.ά. συνδέουν το σύστημα παρακολούθησης μιας νέας γέφυρας (αισθητήρες παραμόρφωσης, στάθμης, επιταχυνσιόμετρα, μετεωρολογικός σταθμός) με μια εφαρμογή AR σε μηχανή παιχνιδιών για το HoloLens. Οι μετρήσεις έρχονται από διαδικτυακή πλατφόρμα και η απεικόνιση γίνεται είτε επιτόπου είτε εξ αποστάσεως \cite{fawad2024bridge}. Οι \en{Awadallah} κ.ά. παρουσιάζουν σύστημα SHM με AR στο HoloLens~2 με έμφαση στην ακριβή ευθυγράμμιση του τρισδιάστατου μοντέλου με την κατασκευή, μέσω αναγνώρισης του μοντέλου από την κάμερα \cite{awadallah2025ar}.

Μια τρίτη ομάδα εργασιών εστιάζει στη διαχείριση της πληροφορίας επιθεώρησης και όχι στις συνεχείς μετρήσεις. Οι \en{Nguyen} κ.ά. αναπτύσσουν την εφαρμογή HoloBridge για το HoloLens. Βασίζεται σε μοντέλο πληροφοριών γέφυρας (BIM) και επιτρέπει την καταγραφή και ενημέρωση ευρημάτων επιθεώρησης και συντήρησης. Δοκιμάστηκε σε υπάρχουσα γέφυρα στη Νότια Κορέα \cite{nguyen2022holobridge}.

\section{BIM και μικτή πραγματικότητα στο εργοτάξιο}
\label{sec:rw-bim}

Η υπέρθεση μοντέλων σχεδιασμού στον πραγματικό χώρο είναι από τις πιο μελετημένες εφαρμογές της MR στις κατασκευές. Οι \en{Chalhoub} και \en{Ayer} μελέτησαν τη χρήση MR για την επικοινωνία του σχεδιασμού ηλεκτρολογικών εγκαταστάσεων σε εργάτες, και κατέγραψαν επιπτώσεις στην απόδοση της εργασίας σε σχέση με τα συμβατικά σχέδια \cite{chalhoub2018mixed}. Στην παρούσα εργασία το μοντέλο (BIM, CAD ή φωτογραμμετρία) δεν είναι ο κύριος σκοπός, αλλά το \emph{πλαίσιο} μέσα στο οποίο διαβάζονται οι μετρήσεις. Τοποθετείται με την ίδια ευθυγράμμιση με τους αισθητήρες, ώστε να μην χρειάζεται δεύτερη διαδικασία.

\todo{Πρόσθεσε 3--5 εργασίες για BIM σε AR/MR στο εργοτάξιο (π.χ. έλεγχος προόδου, σύγκριση as-built/as-designed) και για ψηφιακά δίδυμα υποδομών.}

\section{Γεωαναφορά σε επαυξημένη πραγματικότητα}
\label{sec:rw-georef}

Η τοποθέτηση γεωαναφερμένων γραφικών σε εξωτερικούς χώρους απασχολεί τη βιβλιογραφία εδώ και δύο δεκαετίες. Οι \en{Behzadan} και \en{Kamat} περιέγραψαν μέθοδο ευθυγράμμισης γραφικών κατασκευής σε κινητό σύστημα AR εξωτερικού χώρου, με θέση από GPS και προσανατολισμό από αισθητήρες κεφαλής \cite{behzadan2007georeferenced}. Οι \en{Schall} κ.ά. ανέπτυξαν σύστημα AR χειρός για την απεικόνιση υπόγειων δικτύων υποδομής από δεδομένα GIS, με εντοπισμό RTK-GPS \cite{schall2009handheld}. Πιο πρόσφατα, οι \en{Messi} κ.ά. συνδυάζουν πολλαπλούς μηχανισμούς ευθυγράμμισης χωρίς δείκτες (GNSS-RTK και όραση υπολογιστή) για εφαρμογές διαχείρισης εγκαταστάσεων σε εσωτερικούς και εξωτερικούς χώρους. Αξιολογούν ποσοτικά την ακρίβεια και δείχνουν ότι απαιτείται συχνή επαναευθυγράμμιση για ακρίβεια της τάξης του δεκατόμετρου \cite{messi2024seamless}.

Τα συστήματα αυτά βασίζονται σε GNSS, που δίνει ακρίβεια εκατοστών μόνο με εξοπλισμό RTK, δεν λειτουργεί σε κλειστούς χώρους, σήραγγες ή κάτω από καταστρώματα γεφυρών, και προσθέτει υλικό που πρέπει να συνδεθεί με τη συσκευή. Η παρούσα εργασία ακολουθεί την προσέγγιση των \emph{σημείων ελέγχου} της τοπογραφίας. Ο χρήστης ακουμπά με το χειριστήριο σημεία γνωστών συντεταγμένων, συνήθως τους ίδιους τους αισθητήρες, και το σύστημα λύνει τον μετασχηματισμό. Δεν χρειάζεται επιπλέον εξοπλισμός, λειτουργεί σε κλειστούς χώρους, και η ακρίβεια μπορεί να προβλεφθεί και να ελεγχθεί.

\section{Οπτικοποίηση δεδομένων στη θέση τους}
\label{sec:rw-situated}

Η απεικόνιση δεδομένων στον χώρο όπου αυτά έχουν νόημα μελετάται στην περιοχή της οπτικοποίησης ως \en{situated visualization} και \en{situated analytics}. Οι \en{Willett} κ.ά. διακρίνουν τις \emph{situated} αναπαραστάσεις, που απλώς βρίσκονται κοντά στο φυσικό αντικείμενο, από τις \emph{embedded}, που ενσωματώνονται σε αυτό \cite{willett2017embedded}. Οι \en{ElSayed} κ.ά. παρουσίασαν εργαλεία \en{situated analytics} με AR \cite{elsayed2016situated}. Δύο πρόσφατες επισκοπήσεις οργανώνουν το πεδίο: οι \en{Bressa} κ.ά. εξετάζουν πώς χρησιμοποιείται η έννοια της «θέσης» (χώρος, χρόνος, τόπος, δραστηριότητα, κοινότητα) \cite{bressa2022situated}, και οι \en{Shin} κ.ά. ταξινομούν συστήματα \en{situated analytics} σε τρεις διαστάσεις \cite{shin2023situated}. Ευρύτερα, το πεδίο του \en{immersive analytics} μελετά την ανάλυση δεδομένων σε εμβυθιστικά περιβάλλοντα \cite{marriott2018immersive}.

Με την ορολογία των \en{Willett} κ.ά., τα εικονίδια των αισθητήρων και οι ράβδοι ορίων της παρούσας εργασίας είναι \emph{situated}, ενώ ο χάρτης θερμότητας πάνω στο μοντέλο και τα διανύσματα μετακίνησης πάνω στους στόχους πλησιάζουν τις \emph{embedded} αναπαραστάσεις.

\review{Επιβεβαίωσε ότι η κατάταξη situated/embedded που κάνω εδώ συμφωνεί με τους ορισμούς του \cite{willett2017embedded}. Είναι καλό σημείο για συζήτηση στην εξέταση, άρα πρέπει να το καταλαβαίνεις.}

\section{Σύγκριση}
\label{sec:rw-comparison}

Ο \cref{tab:rw-comparison} συγκρίνει τις εργασίες που βρίσκονται πιο κοντά στο προτεινόμενο σύστημα.

\begin{table}[htbp]
  \centering
  \footnotesize
  \caption{Σύγκριση σχετικών εργασιών.}
  \label{tab:rw-comparison}
  \setlength{\tabcolsep}{4pt}
  \begin{tabularx}{\textwidth}{@{}>{\raggedright\arraybackslash}X
      >{\raggedright\arraybackslash}p{1.6cm} >{\raggedright\arraybackslash}p{1.9cm}
      >{\raggedright\arraybackslash}p{2.6cm} >{\raggedright\arraybackslash}p{1.7cm}
      >{\raggedright\arraybackslash}p{1.6cm}@{}}
    \toprule
    Εργασία & Συσκευή & Ζωντανά δεδομένα αισθητήρων & Ευθυγράμμιση & Μοντέλο 3Δ & Ποσοτική αξιολόγηση ακρίβειας \\
    \midrule
    \cite{fawad2024bridge} & HoloLens & Ναι (web) & \fillin{;} & BIM & \fillin{;} \\
    \cite{awadallah2025ar} & HoloLens 2 & \fillin{;} & Αναγνώριση μοντέλου από κάμερα & Ναι & Ναι \\
    \cite{nguyen2022holobridge} & HoloLens & Όχι (επιθεώρηση) & \fillin{;} & BIM & \fillin{;} \\
    \cite{schall2009handheld} & Χειρός & Όχι (GIS) & RTK-GPS & Δίκτυα & \fillin{;} \\
    \cite{messi2024seamless} & \fillin{;} & Όχι (FM) & GNSS-RTK + όραση & BIM & Ναι \\
    \midrule
    \textbf{Παρούσα εργασία} & Quest 3 & Ναι (REST, polling, push) & Σημεία ελέγχου + spatial anchors & OBJ / DXF / IFC & Ναι \\
    \bottomrule
  \end{tabularx}
\end{table}

\review{Συμπλήρωσε τα \fillin{;} \textbf{μόνο αφού διαβάσεις} τις εργασίες. Αν μια εργασία δεν καλύπτει μια στήλη, γράψε «Όχι» ή «—».}

Από τη σύγκριση προκύπτει ότι οι εργασίες που φέρνουν δεδομένα αισθητήρων στο πεδίο σπάνια αξιολογούν ποσοτικά την ακρίβεια τοποθέτησης. Αντίστροφα, οι εργασίες που εστιάζουν στην ευθυγράμμιση δεν αφορούν ζωντανά δεδομένα παρακολούθησης. Η παρούσα εργασία συνδυάζει: ζωντανά δεδομένα από πλατφόρμα παρακολούθησης, ευθυγράμμιση χωρίς επιπλέον εξοπλισμό με πρόβλεψη σφάλματος, χαμηλού κόστους αυτόνομη συσκευή, και ποσοτική αξιολόγηση στο πεδίο.

\review[!]{Η τελευταία παράγραφος είναι \textbf{ο ισχυρισμός πρωτοτυπίας} της εργασίας. Ισχύει μόνο αν τον στηρίζει η βιβλιογραφία που θα διαβάσεις. Αν βρεις εργασία που κάνει ήδη τον ίδιο συνδυασμό, διατύπωσέ τον πιο μετριοπαθώς.}

\chapter{Απαιτήσεις και σχεδίαση συστήματος}
\label{ch:design}

\guide{Έκταση 12--15 σελίδες. Καλό σημείο εκκίνησης για τη γραφή: το σύστημα το ξέρεις ήδη. Περιγράφεις \emph{τι} και \emph{γιατί}, όχι \emph{πώς} (αυτό είναι το \cref{ch:implementation}).}

\section{Σενάρια χρήσης}
\label{sec:usecases}

\guide{Τρία σενάρια, το καθένα σε μία παράγραφο με τα βήματα του χρήστη:
(α) \textbf{Επιθεώρηση πεδίου}: φόρτωση έργου από Terraweb, ευθυγράμμιση σε 2--4 σημεία ελέγχου, περιήγηση στους αισθητήρες, συναγερμοί, διαδρομή επιθεώρησης, αναφορά.
(β) \textbf{Προετοιμασία στο γραφείο}: το μοντέλο ως μακέτα πάνω στο τραπέζι, χρονομηχανή, χάρτης θερμότητας.
(γ) \textbf{Εκπαίδευση}: το εικονικό έργο εξάσκησης (\code{PracticeProject}) και ο οδηγός πρώτης χρήσης (\code{Onboarding}).}

\section{Απαιτήσεις}
\label{sec:requirements}

\subsection{Λειτουργικές απαιτήσεις}

\begin{table}[htbp]
  \centering\small
  \caption{Λειτουργικές απαιτήσεις.}
  \label{tab:functional-req}
  \begin{tabularx}{\textwidth}{@{}lX@{}}
    \toprule
    Κωδ. & Απαίτηση \\
    \midrule
    ΛΑ1 & Λήψη έργου (σημεία αισθητήρων με συντεταγμένες κανάβου) από την πλατφόρμα παρακολούθησης ή από αρχείο CSV. \\
    ΛΑ2 & Ευθυγράμμιση του έργου με τον πραγματικό χώρο με άγγιγμα γνωστών σημείων και επιβεβαίωση από τον χρήστη. \\
    ΛΑ3 & Διατήρηση της ευθυγράμμισης μεταξύ συνεδριών. \\
    ΛΑ4 & Απεικόνιση κάθε αισθητήρα στη θέση του με τύπο, όνομα, κατάσταση και τελευταία τιμή. \\
    ΛΑ5 & Απεικόνιση ιστορικού μετρήσεων και κατάστασης σε παρελθούσα ημερομηνία. \\
    ΛΑ6 & Ειδοποίηση για αλλαγές κατάστασης συναγερμού σε (σχεδόν) πραγματικό χρόνο. \\
    ΛΑ7 & Υπέρθεση μοντέλου 3Δ (OBJ, DXF, IFC) στο ίδιο σύστημα συντεταγμένων. \\
    ΛΑ8 & Καθοδηγούμενη διαδρομή επιθεώρησης με αναφορά. \\
    ΛΑ9 & Πρόσβαση τρίτων εφαρμογών μέσω τεκμηριωμένου API. \\
    \bottomrule
  \end{tabularx}
\end{table}

\subsection{Μη λειτουργικές απαιτήσεις}

\guide{ΜΛΑ: σταθερός ρυθμός καρέ (72\,Hz), λειτουργία χωρίς σύνδεση (\code{OfflineCache}), ασφάλεια σε τοπικό δίκτυο (κλειδί API, όρια μεγέθους, rate limiting), χρήση με λεπτά γάντια εργασίας (μεγάλοι στόχοι, κουμπιά αντί για χειρονομίες), δίγλωσση διεπαφή (ελληνικά/αγγλικά), ανεξαρτησία από συγκεκριμένο CRS, άνεση (αποφυγή ναυτίας στην κίνηση με stick). Για καθεμία: γιατί προκύπτει από το πεδίο.}

\section{Αρχιτεκτονική}
\label{sec:architecture}

Το σύστημα αποτελείται από τρία επίπεδα (\cref{fig:architecture}): την πλατφόρμα παρακολούθησης στο νέφος, προαιρετικά εργαλεία σε υπολογιστή στο τοπικό δίκτυο, και την εφαρμογή στη συσκευή, η οποία λειτουργεί και αυτόνομα.

\begin{figure}[htbp]
  \centering
  \resizebox{\textwidth}{!}{% Διάγραμμα αρχιτεκτονικής (TikZ). Χρήση: % Διάγραμμα αρχιτεκτονικής (TikZ). Χρήση: \input{../figures/architecture} μέσα σε figure.
% Τα ονόματα κλάσεων/αρχείων είναι τα πραγματικά του QuestMarkerApp.
\begin{tikzpicture}[
  font=\sffamily\scriptsize,
  node distance=4mm and 6mm,
  comp/.style={draw, rounded corners=2pt, fill=white, align=center, minimum width=27mm, minimum height=7mm, inner sep=2pt},
  tier/.style={draw=#1!60!black, fill=#1!8, rounded corners=4pt, inner sep=3mm},
  tiertitle/.style={font=\sffamily\small\bfseries, anchor=south west},
  flow/.style={-{Stealth[length=2mm]}, thick},
  lbl/.style={font=\sffamily\tiny, fill=white, inner sep=1pt, align=center},
]

% ---------------- Cloud -------------------------------------------------
\node[comp] (tw) {Terraweb\\(REST, JWT)};
\node[comp, right=of tw] (wit) {Wit.ai\\(φωνή, προαιρετικό)};

% ---------------- PC ----------------------------------------------------
\node[comp, below=14mm of tw, xshift=-16mm] (bridge) {\texttt{terraweb\_bridge.py}\\\texttt{send\_project.py}};
\node[comp, right=of bridge] (models) {\texttt{send\_model.py}\\\texttt{model\_prep.py}};
\node[comp, right=of models] (cserver) {\texttt{coordinate\_server.py}\\(long-poll \texttt{/api/events})};

% ---------------- Quest: network ------------------------------------------
\node[comp, below=16mm of bridge] (api) {\texttt{LocalAPIServer}\\HTTP :8080};
\node[comp, right=of api] (twc) {\texttt{TerrawebClient}\\\texttt{LiveAlerts}};
\node[comp, right=of twc] (push) {\texttt{WebSocketPushClient}\\\texttt{NetworkDiscovery}};
\node[comp, below=3mm of api, minimum width=27mm] (sec) {\texttt{ApiSecurity}\\κλειδί, rate limit};

% ---------------- Quest: core ---------------------------------------------
\node[comp, below=8mm of sec] (align) {\texttt{ProjectAlignment}\\\texttt{CoordinateTransform}};
\node[comp, right=of align] (anch) {\texttt{SpatialAnchorManager}\\\texttt{AnchorFieldTest}};
\node[comp, right=of anch] (sess) {\texttt{SessionManager}\\\texttt{OfflineCache}};

% ---------------- Quest: visualization / UI -------------------------------
\node[comp, below=of align] (mark) {\texttt{MarkerManager}\\\texttt{HistoryRibbons}};
\node[comp, right=of mark] (model) {\texttt{ModelOverlay}\\\texttt{ModelHeatmap}};
\node[comp, right=of model] (ui) {\texttt{MainMenu}, \texttt{SensorCard}\\\texttt{VoiceCommands}, \texttt{InspectionRoute}};

% ---------------- Tiers ---------------------------------------------------
\begin{scope}[on background layer]
  \node[tier=gray, fit=(tw)(wit)] (cloudT) {};
  \node[tier=green, fit=(bridge)(models)(cserver)] (pcT) {};
  \node[tier=blue, fit=(api)(push)(sec)(align)(sess)(mark)(ui)] (questT) {};
\end{scope}
\node[tiertitle] at (cloudT.north west) {Νέφος};
\node[tiertitle] at (pcT.north west) {PC στο LAN (Python)};
\node[tiertitle] at (questT.north west) {Meta Quest 3 (Unity, C\#)};

% ---------------- Flows ---------------------------------------------------
\draw[flow] (tw) -- node[lbl]{αισθητήρες} (bridge);
\draw[flow] (bridge) -- node[lbl]{POST /api/project} (api);
\draw[flow] (models) -- node[lbl, pos=0.35]{POST /api/model\\(zip)} (api);
\draw[flow] (cserver) -- node[lbl]{push} (push);
\draw[flow, <->] (tw.east) -| ++(4mm,-20mm) -| node[lbl, pos=0.75]{HTTPS: μετρήσεις,\\ιστορικό, συναγερμοί} (twc.north);
\draw[flow, <->, dashed] (wit.south) |- ($(ui.north east)+(3mm,8mm)$) -- ($(ui.north east)+(3mm,0)$);
\draw[flow] (api) -- (sec);
\draw[flow] (sec) -- (align);
\draw[flow] (align) -- (mark);
\draw[flow] (align) -- (anch);
\draw[flow] (align.south east) -- (model);
\draw[flow] (twc) -- ++(0,-8mm) -| (sess);
\end{tikzpicture}
 μέσα σε figure.
% Τα ονόματα κλάσεων/αρχείων είναι τα πραγματικά του QuestMarkerApp.
\begin{tikzpicture}[
  font=\sffamily\scriptsize,
  node distance=4mm and 6mm,
  comp/.style={draw, rounded corners=2pt, fill=white, align=center, minimum width=27mm, minimum height=7mm, inner sep=2pt},
  tier/.style={draw=#1!60!black, fill=#1!8, rounded corners=4pt, inner sep=3mm},
  tiertitle/.style={font=\sffamily\small\bfseries, anchor=south west},
  flow/.style={-{Stealth[length=2mm]}, thick},
  lbl/.style={font=\sffamily\tiny, fill=white, inner sep=1pt, align=center},
]

% ---------------- Cloud -------------------------------------------------
\node[comp] (tw) {Terraweb\\(REST, JWT)};
\node[comp, right=of tw] (wit) {Wit.ai\\(φωνή, προαιρετικό)};

% ---------------- PC ----------------------------------------------------
\node[comp, below=14mm of tw, xshift=-16mm] (bridge) {\texttt{terraweb\_bridge.py}\\\texttt{send\_project.py}};
\node[comp, right=of bridge] (models) {\texttt{send\_model.py}\\\texttt{model\_prep.py}};
\node[comp, right=of models] (cserver) {\texttt{coordinate\_server.py}\\(long-poll \texttt{/api/events})};

% ---------------- Quest: network ------------------------------------------
\node[comp, below=16mm of bridge] (api) {\texttt{LocalAPIServer}\\HTTP :8080};
\node[comp, right=of api] (twc) {\texttt{TerrawebClient}\\\texttt{LiveAlerts}};
\node[comp, right=of twc] (push) {\texttt{WebSocketPushClient}\\\texttt{NetworkDiscovery}};
\node[comp, below=3mm of api, minimum width=27mm] (sec) {\texttt{ApiSecurity}\\κλειδί, rate limit};

% ---------------- Quest: core ---------------------------------------------
\node[comp, below=8mm of sec] (align) {\texttt{ProjectAlignment}\\\texttt{CoordinateTransform}};
\node[comp, right=of align] (anch) {\texttt{SpatialAnchorManager}\\\texttt{AnchorFieldTest}};
\node[comp, right=of anch] (sess) {\texttt{SessionManager}\\\texttt{OfflineCache}};

% ---------------- Quest: visualization / UI -------------------------------
\node[comp, below=of align] (mark) {\texttt{MarkerManager}\\\texttt{HistoryRibbons}};
\node[comp, right=of mark] (model) {\texttt{ModelOverlay}\\\texttt{ModelHeatmap}};
\node[comp, right=of model] (ui) {\texttt{MainMenu}, \texttt{SensorCard}\\\texttt{VoiceCommands}, \texttt{InspectionRoute}};

% ---------------- Tiers ---------------------------------------------------
\begin{scope}[on background layer]
  \node[tier=gray, fit=(tw)(wit)] (cloudT) {};
  \node[tier=green, fit=(bridge)(models)(cserver)] (pcT) {};
  \node[tier=blue, fit=(api)(push)(sec)(align)(sess)(mark)(ui)] (questT) {};
\end{scope}
\node[tiertitle] at (cloudT.north west) {Νέφος};
\node[tiertitle] at (pcT.north west) {PC στο LAN (Python)};
\node[tiertitle] at (questT.north west) {Meta Quest 3 (Unity, C\#)};

% ---------------- Flows ---------------------------------------------------
\draw[flow] (tw) -- node[lbl]{αισθητήρες} (bridge);
\draw[flow] (bridge) -- node[lbl]{POST /api/project} (api);
\draw[flow] (models) -- node[lbl, pos=0.35]{POST /api/model\\(zip)} (api);
\draw[flow] (cserver) -- node[lbl]{push} (push);
\draw[flow, <->] (tw.east) -| ++(4mm,-20mm) -| node[lbl, pos=0.75]{HTTPS: μετρήσεις,\\ιστορικό, συναγερμοί} (twc.north);
\draw[flow, <->, dashed] (wit.south) |- ($(ui.north east)+(3mm,8mm)$) -- ($(ui.north east)+(3mm,0)$);
\draw[flow] (api) -- (sec);
\draw[flow] (sec) -- (align);
\draw[flow] (align) -- (mark);
\draw[flow] (align) -- (anch);
\draw[flow] (align.south east) -- (model);
\draw[flow] (twc) -- ++(0,-8mm) -| (sess);
\end{tikzpicture}
}
  \caption{Αρχιτεκτονική του συστήματος.}
  \label{fig:architecture}
\end{figure}

\guide{Περιέγραψε κάθε επίπεδο και κάθε ροή του σχήματος. Τόνισε την επιλογή: ο \emph{εξυπηρετητής} τρέχει στα γυαλιά (αντί η εφαρμογή να είναι μόνο πελάτης), ώστε οποιαδήποτε συσκευή στο LAN να στέλνει δεδομένα χωρίς ενδιάμεσο. Και ότι η εφαρμογή μιλά και απευθείας με το Terraweb, χωρίς PC.}

\section{Επικοινωνία}
\label{sec:communication}

\subsection{Τοπικό REST API}

\Cref{tab:api-summary} συνοψίζει τα endpoints του ενσωματωμένου εξυπηρετητή (θύρα 8080). Η πλήρης αναφορά βρίσκεται στο \cref{app:api}.

\begin{table}[htbp]
  \centering\small
  \caption{Ομάδες endpoints του τοπικού API.}
  \label{tab:api-summary}
  \begin{tabularx}{\textwidth}{@{}lX@{}}
    \toprule
    Ομάδα & Endpoints \\
    \midrule
    Σημεία & \code{/api/coordinates}, \code{/api/clear}, \code{/api/undo}, \code{/api/layer} \\
    Έργο / ευθυγράμμιση & \code{/api/project}, \code{/api/project/reset}, \code{/api/calibrate/*}, \code{/api/gps-origin}, \code{/api/anchors} \\
    Μοντέλα & \code{/api/model}, \code{/api/models} \\
    Συνεδρίες & \code{/api/session}, \code{/api/session/new} \\
    Πεδίο & \code{/api/fieldnotes}, \code{/api/reports}, \code{/api/fieldtest} \\
    Διάγνωση & \code{/api/health}, \code{/api/status}, \code{/api/logs} \\
    \bottomrule
  \end{tabularx}
\end{table}

\subsection{Ανακάλυψη και push}
\guide{\code{NetworkDiscovery}: εκπομπή UDP στη θύρα 48080 με JSON \{service, port, name\}, ώστε τα εργαλεία PC να βρίσκουν τα γυαλιά χωρίς IP. Push μέσω long-poll στο \code{/api/events} του \code{coordinate_server.py} (\code{WebSocketPushClient}). Γιατί long-poll και όχι WebSocket/MQTT (απλότητα, διέλευση από firewalls, καμία εξάρτηση στη συσκευή).}

\subsection{Σύνδεση με την πλατφόρμα Terraweb}
\guide{Αυθεντικοποίηση JWT (access 15 λεπτά, refresh 7 ημέρες), επιλογή έργου, λήψη αισθητήρων, ιστορικού και ενεργών συναγερμών (polling ανά 20\,s). Μόνο HTTPS για διαπιστευτήρια· pinning πιστοποιητικού για αυτο-υπογεγραμμένους εξυπηρετητές.}

\section{Ασφάλεια}
\label{sec:security}
\guide{Μοντέλο απειλών για LAN: κλειδί API (\code{X-Api-Key} ή \code{Bearer}), όριο ρυθμού αιτημάτων, όρια μεγέθους σώματος, προστασία από zip bomb στα πακέτα μοντέλων, καμία έκθεση εκτός LAN. Τι \emph{δεν} καλύπτεται (cleartext HTTP εντός LAN) και γιατί είναι αποδεκτό εδώ.}

\section{Μοντέλο δεδομένων}
\label{sec:datamodel}
\guide{\code{CoordinateData} (σημείο: id, x/y/z, χρώμα, ετικέτα, στρώμα), \code{ProjectData} (σημεία έργου με id, όνομα, τύπο, E/N/H), συνεδρία, εγγραφές ιστορικού. Σχήμα κλάσεων UML με τα κύρια πεδία.}

\chapter{Υλοποίηση}
\label{ch:implementation}

\guide{Έκταση 20--25 σελίδες. Μία ενότητα ανά υποσύστημα, με το ίδιο μοτίβο: \textbf{πρόβλημα → σχεδιαστική απόφαση και γιατί → (σύντομος) κώδικας ή ψευδοκώδικας → στιγμιότυπο από τα γυαλιά}. \emph{Όχι} περιγραφή κώδικα γραμμή-γραμμή. Αποσπάσματα κώδικα έως 15--20 γραμμές· τα μεγαλύτερα σε ψευδοκώδικα.}

\section{Τεχνολογίες και δομή του έργου}
\label{sec:impl-stack}
\guide{Unity 6 (6000.5), C\#, Meta XR SDK, Built-in Render Pipeline με δικούς shaders, Android/IL2CPP/ARM64. Python 3.11 για τα εργαλεία PC (Flask, requests, pymeshlab, ifcopenshell). Πίνακας φακέλων \code{Assets/Scripts/} (API, Core, Markers, Models, UI, Mobile, Editor) με μία γραμμή περιγραφής ο καθένας. Λειτουργία προεπισκόπησης σε κινητό για ανάπτυξη χωρίς headset.}

% ---------------------------------------------------------------------
\section{Ευθυγράμμιση έργου}
\label{sec:impl-alignment}

Ο χρήστης ακουμπά με την άκρη του χειριστηρίου σημεία ελέγχου με γνωστές συντεταγμένες έργου. Για κάθε ζεύγος $(\mathbf{p}_i, \mathbf{m}_i)$, όπου $\mathbf{p}_i$ η θέση στο τοπικό σύστημα του έργου (\cref{eq:enh-to-local}) και $\mathbf{m}_i$ η μετρημένη θέση στον χώρο της συσκευής, το σύστημα εκτιμά μετασχηματισμό στερεού σώματος τεσσάρων βαθμών ελευθερίας: στροφή $\theta$ γύρω από τον κατακόρυφο άξονα και μετάθεση $\mathbf{t} \in \mathbb{R}^3$. Η κλίση δεν εκτιμάται, γιατί η συσκευή γνωρίζει ήδη την κατακόρυφο από το αδρανειακό της σύστημα με ακρίβεια καλύτερη από όση θα έδιναν λίγα σημεία.

\subsection{Λύση ελαχίστων τετραγώνων}

Με $\bar{\mathbf{p}}, \bar{\mathbf{m}}$ τα κέντρα βάρους και $\mathbf{p}'_i = \mathbf{p}_i - \bar{\mathbf{p}}$, $\mathbf{m}'_i = \mathbf{m}_i - \bar{\mathbf{m}}$, η γωνία που ελαχιστοποιεί το άθροισμα τετραγώνων των οριζόντιων υπολοίπων είναι
\begin{equation}
  \theta = \operatorname{atan2}\!\Big( \sum_i \big(m'_{i,x} p'_{i,z} - m'_{i,z} p'_{i,x}\big),\;
                                      \sum_i \big(m'_{i,x} p'_{i,x} + m'_{i,z} p'_{i,z}\big) \Big),
  \label{eq:yaw}
\end{equation}
και η μετάθεση
\begin{equation}
  \mathbf{t} = \bar{\mathbf{m}} - R_y(\theta)\,\bar{\mathbf{p}} .
  \label{eq:translation}
\end{equation}
Πρόκειται για την ειδική περίπτωση στο επίπεδο της κλειστής λύσης του απόλυτου προσανατολισμού \cite{horn1987closed,arun1987leastsquares}. Η κατακόρυφη συνιστώσα της $\mathbf{t}$ είναι η μέση κατακόρυφη απόκλιση.

\guide{Εξήγησε την παραγωγή της \cref{eq:yaw} σε 4--5 γραμμές (μεγιστοποίηση του $\sum \mathbf{m}'_i \cdot R\mathbf{p}'_i$). Υλοποίηση: \code{ProjectAlignment.Solve} στο \code{Assets/Scripts/Core/ProjectAlignment.cs}.}

\subsection{Ποιότητα και πρόβλεψη σφάλματος}

Μετά από κάθε άγγιγμα υπολογίζονται το οριζόντιο και το κατακόρυφο RMS των υπολοίπων, το μέγιστο υπόλοιπο και το πιθανότερο λανθασμένο σημείο. Επιπλέον, το σύστημα \emph{προβλέπει} το πλευρικό σφάλμα στο πιο απομακρυσμένο σημείο του έργου. Αν κάθε άγγιγμα έχει τυπικό σφάλμα $\sigma$ (στην υλοποίηση $\sigma = \SI{2}{cm}$), η γωνιακή αβεβαιότητα από βάση μήκους $b$ είναι περίπου $\sigma\sqrt{2}/b$, οπότε σε απόσταση $D$ από το κέντρο των σημείων ελέγχου:
\begin{equation}
  \varepsilon_{\text{πλευρ.}}(D) \approx \sigma \sqrt{2}\, \frac{D}{b} .
  \label{eq:predicted-error}
\end{equation}
Η \cref{eq:predicted-error} εμφανίζεται στον χρήστη και ελέγχεται πειραματικά στην \cref{sec:exp-alignment}.

\subsection{Καθοδήγηση του χρήστη}

\begin{algorithm}[htbp]
  \caption{Πρόταση επόμενου σημείου ελέγχου (farthest point).}
  \label{alg:suggest-next}
  \begin{algorithmic}[1]
    \Require σημεία έργου $P$, ήδη τοποθετημένα $T \subset P$
    \State $best \gets \text{κανένα}$, $score \gets -1$
    \ForAll{$q \in P \setminus T$}
      \State $d \gets \min_{p \in T} \lVert q_{xz} - p_{xz} \rVert$
      \If{$d > score$} $score \gets d$, $best \gets q$ \EndIf
    \EndFor
    \State \Return $best$
  \end{algorithmic}
\end{algorithm}

\guide{Επιπλέον έλεγχοι: σημεία πολύ κοντά για να οριστεί γωνία ($< \SI{25}{cm}$), δύο αγγίγματα στο ίδιο σημείο, ανίχνευση κατοπτρισμού (αντιμετατεθειμένες στήλες E/N: με $\geq 3$ σημεία και RMS $> \SI{10}{cm}$ ελέγχεται αν η κατοπτρική λύση ταιριάζει πολύ καλύτερα), επανέλεγχος αποθηκευμένης ευθυγράμμισης με ένα άγγιγμα (ανοχή \SI{10}{cm}). Στιγμιότυπο του πίνακα ευθυγράμμισης με τα RMS.}

% ---------------------------------------------------------------------
\section{Spatial anchors και επαναφορά}
\label{sec:impl-anchors}
\guide{\code{SpatialAnchorManager}: κύκλος ζωής anchors. \code{AnchorFieldTest} + \code{FieldTestLog}: ένα anchor ανά σημείο ελέγχου, μέτρηση επανεντοπισμού και μετατόπισης, CSV ανά έργο και ημέρα (στήλες στο \cref{app:raw}). Γιατί η επαναφορά δεν γίνεται αυτόματα χωρίς επιβεβαίωση (εμπιστοσύνη πεδίου).}

% ---------------------------------------------------------------------
\section{Απόδοση σημείων (markers)}
\label{sec:impl-markers}
\guide{\code{MarkerManager}/\code{Marker}: billboard sprites με δικό shader (\code{MarkerSprite}), batching για λίγα draw calls, αποφυγή δεσμεύσεων μνήμης ανά καρέ (GC), single-pass instanced. Κλίμακα που μένει αναγνώσιμη από κοντά και από μακριά. Ετικέτες, ομαδοποίηση, απόκρυψη από πραγματικά αντικείμενα με το depth API (\code{DepthOcclusion}).}

% ---------------------------------------------------------------------
\section{Οπτικοποίηση δεδομένων}
\label{sec:impl-vis}

\subsection{Κορδέλες ιστορικού και χρονομηχανή}
\guide{\code{HistoryRibbons}: διάγραμμα 30 ημερών δίπλα σε κάθε αισθητήρα. \code{SensorHistory}: λήψη/cache χρονοσειρών. Χρονομηχανή: όλη η σκηνή σε παρελθούσα ημερομηνία.}

\subsection{Χάρτης θερμότητας μετακινήσεων}
Ο χάρτης θερμότητας (\code{ModelHeatmap}) χρωματίζει την επιφάνεια του μοντέλου με την παρεμβαλλόμενη μετακίνηση. Οι αισθητήρες ενώνονται σε τρίγωνα σε κάτοψη με τριγωνοποίηση Delaunay (\code{Delaunay.cs}) και η τιμή σε κάθε σημείο προκύπτει με γραμμική παρεμβολή μέσα στο τρίγωνο. Τρίγωνα με ακμή μεγαλύτερη από δύο φορές την εμβέλεια απορρίπτονται, ώστε να μη ``βάφεται'' κενό χωρίς αισθητήρες. Η εμβέλεια προκύπτει από την τυπική απόσταση των αισθητήρων $s$:
\begin{equation}
  r = \operatorname{clamp}(1{,}5\,s,\ \SI{5}{m},\ \SI{30}{m}) \cdot k,
\end{equation}
με $k$ τον συντελεστή της ρύθμισης στενό/κανονικό/ευρύ, και το χρώμα σβήνει σταδιακά πέρα από αυτήν.

\guide{Γιατί αντικαταστάθηκε η αρχική IDW (commit \code{fd79f04} → \code{1b0a854}): τα ``μάτια ταύρου'' γύρω από κάθε αισθητήρα. Τα τρία μέτρα (συνολική, κατακόρυφη $dz$, ρυθμός), υπόμνημα, τιμή κάτω από τον δείκτη μέσω πλέγματος υψών του μοντέλου. Σύγκριση των μεθόδων στην \cref{sec:exp-interpolation}.}

\subsection{Συναγερμοί και όρια}
\guide{\code{LiveAlerts} (polling 20\,s, αλλαγή κατάστασης → εικονίδιο, ήχος, μήνυμα· συμπεριφορά χωρίς σύνδεση), \code{AlarmBeacons}, \code{SensorLimits}/\code{LimitBars} (ποσοστό εξάντλησης ορίου).}

% ---------------------------------------------------------------------
\section{Μοντέλα 3Δ}
\label{sec:impl-models}
\guide{Αλυσίδα: OBJ/DXF/IFC → (PC) \code{model_prep.py}: απλοποίηση QEM με MeshLab \cite{garland1997qem,cignoni2008meshlab} σε 400k τρίγωνα, σμίκρυνση υφών σε 2048\,px, προαιρετική μεταφορά στον κάναβο των αισθητήρων (\code{model_to_local.py}, ευθυγράμμιση με κοινά σημεία) → zip → \code{POST /api/model} → (Quest) αποσυμπίεση με όρια ασφαλείας, ανάλυση σε νήμα παρασκηνίου, τεμαχισμός σε meshes με δείκτες 16 bit (\code{ModelOverlay}), shader \code{ModelTextured}. Προβολή: μακέτα / 1:1 / πάνω στο έργο. Κίνηση μέσα στο 1:1 με stick και vignette άνεσης (\code{ComfortVignette}).}

% ---------------------------------------------------------------------
\section{Αλληλεπίδραση}
\label{sec:impl-interaction}
\guide{Ακτίνα και κλικ (\code{VrPointer}), κύριο μενού, κάρτα αισθητήρα, πληροφορία με hover, μετροταινία (\code{MeasureTool}), διαδρομή επιθεώρησης (\code{InspectionRoute}: βέλη στο δάπεδο, τικ, αναφορά), φωνητικές εντολές (\code{VoiceCommands}/\code{VoiceRecognizers}, Wit.ai, εναλλακτικά πληκτρολόγιο VR), οδηγός πρώτης χρήσης, υποδείξεις χειριστηρίων, αριστερόχειρες. Σχεδιαστικές αρχές: λεπτά γάντια, δύο γλώσσες, επιβεβαίωση κρατώντας πατημένο (\code{HoldConfirm}) για καταστροφικές ενέργειες.}

% ---------------------------------------------------------------------
\section{Απόδοση της συσκευής}
\label{sec:impl-performance}
\guide{\code{DisplayPerformance}: 72/90\,Hz, δυναμικό foveated rendering. Μετρήσεις στην \cref{sec:exp-performance}.}

% ---------------------------------------------------------------------
\section{Εργαλεία PC και έλεγχοι}
\label{sec:impl-tools}
\guide{\code{terraweb_bridge.py}, \code{send_project.py}, \code{send_model.py}, \code{coordinate_server.py}, \code{fake_scenarios.py} (εξομοίωση σεναρίων). 127 unit tests με pytest (μετασχηματισμοί, προετοιμασία μοντέλων, API).}

\chapter{Πειραματική αξιολόγηση}
\label{ch:evaluation}

% Εικόνα από το analysis/ αν υπάρχει, αλλιώς πλαίσιο-υπενθύμιση.
\newcommand{\genfig}[2][\linewidth]{%
  \IfFileExists{../figures/generated/#2}%
    {\includegraphics[width=#1]{#2}}%
    {\fbox{\parbox[c][4cm][c]{0.8\linewidth}{\centering\sffamily\small Θα παραχθεί από \code{thesis/analysis/}\\{\scriptsize\code{#2}}}}}}
\newcommand{\gentable}[1]{%
  \IfFileExists{../figures/generated/#1}%
    {\begin{adjustbox}{max width=\linewidth}\input{../figures/generated/#1}\end{adjustbox}}%
    {\fillin{πίνακας \code{#1} από το analysis/}}}

Το κεφάλαιο αξιολογεί το σύστημα ως προς τα τέσσερα ερευνητικά ερωτήματα της \cref{sec:rq}. Για κάθε πείραμα περιγράφονται ο στόχος, η διάταξη, η διαδικασία και τα μέτρα. Ακολουθούν τα αποτελέσματα και η συζήτησή τους. Όλοι οι αριθμοί, πίνακες και γραφήματα παράγονται από τα σενάρια ανάλυσης του φακέλου \code{thesis/analysis/} πάνω στα ακατέργαστα δεδομένα του \code{thesis/data/}, ώστε να μπορούν να αναπαραχθούν.

\review[!]{\textbf{Όλο το κεφάλαιο περιμένει μετρήσεις.} Το κείμενο της μεθοδολογίας είναι γραμμένο ως σχέδιο: όπου η πραγματική διαδικασία διαφέρει (αριθμός επαναλήψεων, αποστάσεις, χώρος), \textbf{διόρθωσε το κείμενο ώστε να περιγράφει ό,τι έκανες πραγματικά}, όχι ό,τι σχεδιάστηκε. Τα αποτελέσματα και η συζήτηση γράφονται από σένα, με βάση τα γραφήματα.}

\section{Πειραματική διάταξη}
\label{sec:exp-setup}

Όλα τα πειράματα έγιναν με Meta Quest~3 (λογισμικό συσκευής \fillin{έκδοση}) και την έκδοση της εφαρμογής που αντιστοιχεί στην ετικέτα \code{thesis-v2} του αποθετηρίου. Τα πειράματα πεδίου έγιναν στο \fillin{έργο / χώρος, περιγραφή, έκταση, αριθμός αισθητήρων}. Οι συντεταγμένες αναφοράς των σημείων προέρχονται από τοπογραφική αποτύπωση με \fillin{όργανο} με εκτιμώμενη ακρίβεια \fillin{x}\,mm. Η σύνδεση έγινε μέσω \fillin{Wi-Fi 5\,GHz / hotspot κινητού}.

\todo{Φωτογραφία ή κάτοψη του χώρου των πειραμάτων με τα σημεία ελέγχου και τα ανεξάρτητα σημεία.}

% ---------------------------------------------------------------------
\section{Ακρίβεια ευθυγράμμισης (ΕΕ1)}
\label{sec:exp-alignment}

\IfFileExists{../figures/generated/alignment_numbers.tex}{\input{../figures/generated/alignment_numbers.tex}}{}

\paragraph{Στόχος.} Να μετρηθεί πόσο απέχει ένα σημείο από τη σωστή του θέση μετά την ευθυγράμμιση, ως συνάρτηση του πλήθους και της γεωμετρίας των σημείων ελέγχου, και να ελεγχθεί πόσο καλά το προβλέπει η \cref{eq:predicted-error}.

\paragraph{Διάταξη.} Επιλέγονται 6--10 σημεία με γνωστές συντεταγμένες, κατανεμημένα στον χώρο σε αποστάσεις έως \fillin{x}\,m. Για κάθε επανάληψη, $k \in \{2, 3, 4\}$ από αυτά χρησιμοποιούνται για την ευθυγράμμιση και τα υπόλοιπα ως \emph{ανεξάρτητα σημεία ελέγχου}, που δεν συμμετέχουν στη λύση. Κάθε συνδυασμός επαναλαμβάνεται \fillin{n} φορές από την αρχή.

\paragraph{Εργαλείο μέτρησης.} Για το πείραμα προστέθηκε στην εφαρμογή η λειτουργία «Σημεία ελέγχου ακρίβειας» (\code{AlignmentAccuracy}). Ανοίγει μαζί με τη δοκιμή anchors από μία ροή, «Ξεκινάω μετρήσεις» (Μενού › Ευθυγράμμιση), που ζητά πρώτα τις συνθήκες και αν άλλαξε κάτι στο εργοτάξιο, και είναι κλειδωμένη στη λειτουργία Γραφείο. Ενεργοποιείται μόνο όταν υπάρχει ευθυγράμμιση, και τότε «παγώνει» την τρέχουσα λύση ως μια \emph{σειρά} μετρήσεων, με το πλήθος $k$ και τα αναγνωριστικά των σημείων ελέγχου, τη βάση $b$ και το RMS της λύσης. Όσο είναι ενεργή, ένα άγγιγμα σημείου του έργου \emph{δεν} τοποθετεί σημείο και δεν αλλάζει τη λύση. Η θέση του αγγίγματος μεταφέρεται στον κάναβο του έργου μέσω της αντίστροφης του μετασχηματισμού (\cref{eq:yaw,eq:translation}) και συγκρίνεται με τις συντεταγμένες του σημείου. Η διαφορά $(\Delta E, \Delta N, \Delta H)$ εμφανίζεται αμέσως στο πάνελ και γράφεται σε αρχείο CSV, μαζί με την απόσταση $D$ του σημείου από το κέντρο των σημείων ελέγχου και το προβλεπόμενο σφάλμα της \cref{eq:predicted-error}. Κάθε γραμμή χαρακτηρίζεται ως \emph{ανεξάρτητη} μέτρηση (σημείο εκτός της λύσης) ή ως μέτρηση \emph{σημείου ελέγχου} (σημείο της λύσης, ξανά-αγγιγμένο). Το δεύτερο μετρά την επαναληψιμότητα του αγγίγματος. Μια λανθασμένη μέτρηση αναιρείται επιτόπου με το κουμπί B, και μια νέα ευθυγράμμιση ξεκινά νέα σειρά. Τα αρχεία διατίθενται μέσω \code{GET /api/accuracy} και αναλύονται με το \code{analysis/alignment_accuracy.py}.

Η μέτρηση γίνεται στον κάναβο του έργου και όχι στον χώρο της συσκευής. Έτσι το σφάλμα αναλύεται σε συνιστώσες Ανατολής, Βορρά και ύψους, όπως σε μια τοπογραφική αποτύπωση, και είναι ανεξάρτητο από τη θέση της αρχής παρακολούθησης της συσκευής.

\paragraph{Διαδικασία.} Για κάθε σειρά: ευθυγράμμιση με $k$ σημεία· ενεργοποίηση της μέτρησης· άγγιγμα κάθε ανεξάρτητου σημείου μία φορά· άγγιγμα ενός σημείου ελέγχου \fillin{3} φορές για την εκτίμηση της επαναληψιμότητας· απενεργοποίηση. Οι συνθήκες φωτισμού δηλώνονται στο μενού και καταγράφονται σε κάθε γραμμή.

\paragraph{Μέτρα.} Το RMS της λύσης στα σημεία ελέγχου, όπως το δίνει η εφαρμογή. Το οριζόντιο σφάλμα $e_H = \sqrt{\Delta E^2 + \Delta N^2}$ και το κατακόρυφο $e_V = \Delta H$ στα ανεξάρτητα σημεία (διάμεσος, 95ο εκατοστημόριο). Η σχέση του $e_H$ με την απόσταση $D$. Ο λόγος μετρημένου προς προβλεπόμενο σφάλμα. Η επαναληψιμότητα του αγγίγματος: από τα επαναλαμβανόμενα αγγίγματα του ίδιου σημείου στην ίδια σειρά υπολογίζεται η συγκεντρωτική τυπική απόκλιση ανά άξονα,
\begin{equation}
  \hat\sigma_E^2 = \frac{\sum_j \sum_i (E_{ij} - \bar E_j)^2}{\sum_j (n_j - 1)},
\end{equation}
όπου $j$ τα σημεία και $i$ τα αγγίγματα, και ομοίως για $N$ και $H$. Η οριζόντια εκτίμηση $\hat\sigma = \sqrt{(\hat\sigma_E^2 + \hat\sigma_N^2)/2}$ συγκρίνεται με την τιμή $\sigma = \SI{2}{cm}$ που υποθέτει η \cref{eq:predicted-error}.

\begin{table}[htbp]
  \centering\small
  \caption{Ακρίβεια ευθυγράμμισης ανά πλήθος σημείων ελέγχου $k$ ($N$: ανεξάρτητες μετρήσεις).}
  \label{tab:alignment-by-k}
  \gentable{alignment_by_k.tex}
\end{table}

\begin{figure}[htbp]
  \centering
  \begin{subfigure}{0.56\linewidth}
    \genfig{alignment_error_vs_distance.pdf}
    \caption{Σφάλμα ως προς την απόσταση $D$. Διακεκομμένες: πρόβλεψη.}
    \label{fig:alignment-error}
  \end{subfigure}\hfill
  \begin{subfigure}{0.42\linewidth}
    \genfig{alignment_pred_vs_meas.pdf}
    \caption{Μετρημένο έναντι προβλεπόμενου.}
    \label{fig:alignment-pred}
  \end{subfigure}
  \caption{Οριζόντιο σφάλμα στα ανεξάρτητα σημεία, για $k = 2, 3, 4$.}
  \label{fig:alignment}
\end{figure}

\paragraph{Αποτελέσματα.}\mbox{}

\IfFileExists{../figures/generated/alignment_numbers.tex}{%
Από \AccChecks{} ανεξάρτητες μετρήσεις, το διάμεσο οριζόντιο σφάλμα ήταν \AccMedianCheckHcm\,cm και το 95ο εκατοστημόριο \AccPNinetyFiveCheckHcm\,cm. Η επαναληψιμότητα του αγγίγματος εκτιμήθηκε σε $\hat\sigma = \AccSigmaHcm$\,cm οριζόντια ανά άξονα και \AccSigmaVcm\,cm κατακόρυφα (\AccSigmaDof{} βαθμοί ελευθερίας).}{}
\todo{Σχόλιο στον \cref{tab:alignment-by-k} και στο \cref{fig:alignment}: αυξάνει το σφάλμα γραμμικά με $D$ και μειώνεται με τη βάση $b$; Υπερεκτιμά ή υποεκτιμά η πρόβλεψη (λόγος μετρ./προβλ.); Συμφωνεί το μετρημένο $\hat\sigma$ με τα \SI{2}{cm}; Τι σημαίνουν τα νούμερα για την πράξη (π.χ. πόσα σημεία, σε ποια απόσταση);}

% ---------------------------------------------------------------------
\section{Σταθερότητα και επανεντοπισμός anchors (ΕΕ2)}
\label{sec:exp-anchors}

\paragraph{Στόχος.} Να μετρηθεί αν τα spatial anchors ανά σημείο ελέγχου (\cref{sec:impl-anchors}) ξαναβρίσκονται σε επόμενες επισκέψεις, πόσο γρήγορα, και πόσο έχουν μετακινηθεί, υπό διαφορετικές συνθήκες.

\paragraph{Διαδικασία.} Τρία σημεία ελέγχου σε αποστάσεις \SIrange{5}{30}{m} μεταξύ τους. Την πρώτη ημέρα η ευθυγράμμιση γίνεται κανονικά με ενεργή τη «Δοκιμή πεδίου anchors», που καρφώνει ένα anchor σε κάθε σημείο. Ακολουθούν επισκέψεις μετά από 10 λεπτά, μία ημέρα και μία εβδομάδα. Σε κάθε επίσκεψη η εφαρμογή φορτώνει και επανεντοπίζει τα anchors (έως \SI{20}{s}) και καταγράφει το αποτέλεσμα. Στη συνέχεια ο χρήστης ακουμπά ξανά τα ίδια σημεία, και καταγράφεται η απόκλιση κάθε anchor από το νέο άγγιγμα. Οι επισκέψεις επαναλαμβάνονται υπό διαφορετικό φωτισμό (ηλιοφάνεια, συννεφιά, πρωί, μεσημέρι, βράδυ), μετά από αλλαγές στο εργοτάξιο, και μετά από απομάκρυνση \SI{50}{m} και επιστροφή χωρίς κλείσιμο της εφαρμογής. Οι συνθήκες δηλώνονται από τον χρήστη πριν από κάθε μέτρηση και καταγράφονται σε κάθε γραμμή.

\paragraph{Μέτρα.} Ποσοστό επιτυχούς επανεντοπισμού και διάμεσος χρόνος (γραμμές \code{load}/\code{manual}). Οριζόντια απόκλιση anchor--νέου αγγίγματος (διάμεσος, 95ο εκατοστημόριο) και απόλυτη κατακόρυφη απόκλιση (γραμμές \code{touch}). Τα μέτρα δίνονται ανά συνθήκη και ως συνάρτηση της ηλικίας του anchor. Η ανάλυση γίνεται με το \code{analysis/fieldtest_anchors.py}.

\begin{table}[htbp]
  \centering\small
  \caption{Επανεντοπισμός και απόκλιση anchors ανά συνθήκη.}
  \label{tab:anchors-by-condition}
  \gentable{anchors_by_condition.tex}
\end{table}

\begin{figure}[htbp]
  \centering
  \begin{subfigure}{0.49\linewidth}
    \genfig{anchors_offset_by_condition.pdf}
    \caption{Οριζόντια απόκλιση ανά συνθήκη.}
  \end{subfigure}\hfill
  \begin{subfigure}{0.49\linewidth}
    \genfig{anchors_offset_vs_age.pdf}
    \caption{Απόκλιση συναρτήσει ηλικίας anchor.}
  \end{subfigure}
  \caption{Σταθερότητα των spatial anchors.}
  \label{fig:anchors}
\end{figure}

\paragraph{Αποτελέσματα.}\mbox{}

\todo{Σχόλιο στον \cref{tab:anchors-by-condition} και στο \cref{fig:anchors}: σε ποιες συνθήκες τα anchors είναι αξιόπιστα; Αυξάνεται η απόκλιση με τον χρόνο; Τι σημαίνει αυτό για την αυτόματη επαναευθυγράμμιση;}

% ---------------------------------------------------------------------
\section{Απόδοση γραφικών (ΕΕ3)}
\label{sec:exp-performance}

\paragraph{Στόχος.} Να βρεθούν τα όρια φόρτου για σταθερό ρυθμό καρέ και η επίδραση του ρυθμού ανανέωσης και του foveated rendering.

\paragraph{Διάταξη.} Το εργαλείο OVR Metrics Tool καταγράφει τον ρυθμό καρέ, τους χρόνους CPU και GPU, τη θερμοκρασία και την μπαταρία σε αρχείο CSV. Τα σενάρια φόρτου είναι σωρευτικά:
(α) 50, 200 και 1000 σημεία·
(β) με ετικέτες, κορδέλες ιστορικού και ράβδους ορίων·
(γ) με μοντέλο 100.000, 300.000, 500.000 και 1.000.000 τριγώνων·
(δ) με χάρτη θερμότητας.
Κάθε σενάριο μετριέται για 2 λεπτά σε τρεις ρυθμίσεις: \SI{72}{Hz} χωρίς foveation, \SI{90}{Hz} χωρίς foveation, \SI{90}{Hz} με υψηλό foveation. Τα σημεία δημιουργούνται με το \code{fake_scenarios.py}, ώστε το σενάριο να είναι ίδιο σε κάθε επανάληψη.

\paragraph{Μέτρα.} Μέσος ρυθμός καρέ και 1\% χαμηλότερος, χρόνος καρέ GPU και CPU, ποσοστό καρέ που χάθηκαν, θερμοκρασία, κατανάλωση μπαταρίας ανά 10 λεπτά. Επιπλέον, ο χρόνος φόρτωσης μοντέλου (λήψη, ανάλυση, κατασκευή meshes) ως συνάρτηση του μεγέθους του αρχείου.

\begin{figure}[htbp]
  \centering
  \genfig[0.8\linewidth]{fps_vs_triangles.pdf}
  \caption{Ρυθμός καρέ συναρτήσει του πλήθους τριγώνων του μοντέλου, για τις τρεις ρυθμίσεις.}
  \label{fig:fps-triangles}
\end{figure}

\paragraph{Αποτελέσματα.}\mbox{}

\todo{Πού πέφτει ο ρυθμός κάτω από τα 72/90; Πόσο βοηθά το foveation; Επιβεβαιώνεται το όριο των 400.000 τριγώνων που χρησιμοποιεί η προετοιμασία μοντέλων;}

% ---------------------------------------------------------------------
\section{Προετοιμασία μοντέλων φωτογραμμετρίας}
\label{sec:exp-modelprep}

\paragraph{Στόχος.} Να ποσοτικοποιηθεί ο συμβιβασμός ποιότητας--απόδοσης της απλοποίησης (\cref{sec:impl-models}).

\paragraph{Διαδικασία.} Ένα μοντέλο φωτογραμμετρίας της κατασκευής απλοποιείται με το \code{model_prep.py} σε 1.000.000, 400.000, 200.000 και 100.000 τρίγωνα. Για κάθε εκδοχή μετρώνται: η απόσταση Hausdorff (μέση και μέγιστη) από το αρχικό πλέγμα με το pymeshlab \cite{cignoni1998metro}, το μέγεθος του πακέτου, ο χρόνος μεταφοράς και φόρτωσης, και ο ρυθμός καρέ (σενάριο (γ) της \cref{sec:exp-performance}).

\paragraph{Αποτελέσματα.}\mbox{}

\todo{Γράφημα: γεωμετρικό σφάλμα και ρυθμός καρέ συναρτήσει τριγώνων. Δικαιολογεί η καμπύλη την προεπιλογή των 400.000;}

% ---------------------------------------------------------------------
\section{Καθυστέρηση από άκρο σε άκρο}
\label{sec:exp-latency}

\paragraph{Στόχος.} Να μετρηθεί ο χρόνος από την αποστολή ενός σημείου ως την εμφάνισή του στη συσκευή.

\paragraph{Διαδικασία.} Ο εξυπηρετητής της συσκευής απαντά στο \code{POST /api/coordinates} μόλις βάλει τα σημεία στην ουρά του κύριου νήματος, πριν αυτά εμφανιστούν. Ο χρόνος απόκρισης του POST μετρά επομένως μόνο τη μεταφορά. Για τον χρόνο ως την εμφάνιση, ένα σενάριο Python στέλνει σημεία με μοναδικά αναγνωριστικά (1, 100 και 1000 ανά αίτημα) και αμέσως μετά ρωτά επανειλημμένα το \code{GET /api/coordinates}, που εξυπηρετείται από το κύριο νήμα, έως ότου εμφανιστούν όλα. Ο χρόνος από την αποστολή ως εκείνη την απάντηση είναι ένα άνω όριο της καθυστέρησης από άκρο σε άκρο, με ακρίβεια ενός καρέ (\SI{14}{ms} στα \SI{72}{Hz}). Η ίδια μέτρηση γίνεται και για τη διαδρομή push μέσω του \code{coordinate_server.py}. Δίκτυα: Wi-Fi \SI{5}{GHz} και hotspot κινητού.

\paragraph{Αποτελέσματα.}\mbox{}

\todo{Πίνακας: διάμεσος και 95ο εκατοστημόριο ανά μέγεθος αιτήματος και δίκτυο.}

% ---------------------------------------------------------------------
\section{Σύγκριση μεθόδων παρεμβολής}
\label{sec:exp-interpolation}

\paragraph{Στόχος.} Να συγκριθεί η γραμμική παρεμβολή σε τριγωνοποίηση Delaunay (η μέθοδος της εφαρμογής) με την IDW (η αρχική μέθοδος) στα πραγματικά δεδομένα του έργου.

\paragraph{Διαδικασία.} Διασταυρωμένη επικύρωση «αφήνοντας ένα έξω» (\en{leave-one-out}): για κάθε αισθητήρα, η μετακίνησή του εκτιμάται από τους υπόλοιπους με κάθε μέθοδο, και καταγράφεται η διαφορά από τη μετρημένη τιμή. Η ανάλυση γίνεται εκτός συσκευής, σε Python, με τις μετακινήσεις όλων των τοπογραφικών στόχων σε \fillin{n} ημερομηνίες. Μέθοδοι: γραμμική σε Delaunay, IDW με $p = 1, 2, 3$, και πλησιέστερος γείτονας ως βάση σύγκρισης.

\paragraph{Μέτρα.} RMSE και μέσο απόλυτο σφάλμα ανά μέθοδο, χωριστά για αισθητήρες στο εσωτερικό και στην περίμετρο της διάταξης.

\paragraph{Αποτελέσματα.}\mbox{}

\todo{Πίνακας RMSE/MAE. Επιβεβαιώνει η μέτρηση την επιλογή της γραμμικής παρεμβολής;}

% ---------------------------------------------------------------------
\section{Μελέτη χρηστών (ΕΕ4)}
\label{sec:exp-userstudy}

\paragraph{Στόχος.} Να διαπιστωθεί αν η επιτόπια απεικόνιση βελτιώνει την ταχύτητα και την ακρίβεια εργασιών επιθεώρησης σε σχέση με τη συμβατική ροή, και πώς αξιολογούν οι χρήστες τη χρηστικότητά της.

\paragraph{Συμμετέχοντες.} \fillin{N} συμμετέχοντες (\fillin{μηχανικοί / φοιτητές}), \fillin{x} με προηγούμενη εμπειρία σε VR/AR. Η συμμετοχή ήταν εθελοντική, με γραπτή συναίνεση (\cref{app:userstudy}).

\paragraph{Σχεδιασμός.} Σχεδιασμός εντός υποκειμένων (\en{within-subjects}) με δύο συνθήκες: \textbf{Α}, κάτοψη και λίστα αισθητήρων σε tablet· \textbf{Β}, η εφαρμογή. Η σειρά των συνθηκών αντισταθμίζεται: οι μισοί συμμετέχοντες ξεκινούν από την Α. Για να μην ωφελεί η εξοικείωση τη δεύτερη συνθήκη, χρησιμοποιούνται δύο ισοδύναμα σύνολα αισθητήρων. Πριν από τη συνθήκη Β κάθε συμμετέχων εξασκείται με το έργο εξάσκησης για \fillin{5} λεπτά.

\paragraph{Εργασίες.} (1) Εντοπισμός και άγγιγμα τριών συγκεκριμένων αισθητήρων. (2) Εύρεση των αισθητήρων σε κατάσταση συναγερμού και των θέσεών τους. (3) Εντοπισμός της περιοχής με τη μεγαλύτερη μετακίνηση τον τελευταίο μήνα.

\paragraph{Μέτρα.} Χρόνος ολοκλήρωσης και λάθη ανά εργασία. Χρηστικότητα με το SUS \cite{brooke1996sus} και φόρτος εργασίας με το NASA-TLX \cite{hart1988nasatlx}, και τα δύο μετά από κάθε συνθήκη. Συμπτώματα ναυτίας με το SSQ \cite{kennedy1993ssq} μετά τη συνθήκη Β. Ημιδομημένη συνέντευξη στο τέλος.

\paragraph{Ανάλυση.} Λόγω του μικρού δείγματος και των ζευγαρωτών μετρήσεων, οι συνθήκες συγκρίνονται με το μη παραμετρικό κριτήριο Wilcoxon \cite{wilcoxon1945}, με επίπεδο σημαντικότητας 0{,}05, και αναφέρεται μέγεθος επίδρασης.

\review[!]{Η μελέτη χρηστών χρειάζεται \textbf{έγκριση από την επιτροπή ηθικής} του ιδρύματος, αν το απαιτεί ο κανονισμός, και ρεαλιστικό αριθμό συμμετεχόντων. Αν γίνει πιλοτική μελέτη με 5--6 άτομα, δήλωσέ το ρητά και μην κάνεις στατιστικούς ελέγχους σημαντικότητας: δώσε περιγραφικά στατιστικά.}

\paragraph{Αποτελέσματα.}\mbox{}

\todo{Πίνακας χρόνων/λαθών ανά εργασία και συνθήκη, SUS και NASA-TLX, αποτελέσματα Wilcoxon, κύρια ευρήματα της συνέντευξης.}

% ---------------------------------------------------------------------
\section{Ποιότητα λογισμικού}
\label{sec:exp-quality}

Τα εργαλεία Python ελέγχονται με 127 unit tests (pytest). Καλύπτουν τους μετασχηματισμούς συντεταγμένων, την προετοιμασία μοντέλων και τη μεταφορά τους στον κάναβο των αισθητήρων, την αποστολή μοντέλων, τους πελάτες και το API του εξυπηρετητή push. Όλα περνούν στην έκδοση \code{thesis-v2}. Ο κώδικας C\# της εφαρμογής μεταγλωττίζεται χωρίς σφάλματα και χωρίς προειδοποιήσεις μεταγλωττιστή.

\todo{Προαιρετικά: κάλυψη κώδικα με \code{pytest --cov} και πίνακας ανά αρχείο.}

\chapter{Συζήτηση}
\label{ch:discussion}

\guide{Έκταση 4--6 σελίδες.}

\section{Απαντήσεις στα ερευνητικά ερωτήματα}
\label{sec:answers}
\guide{Μία παράγραφος για καθένα από τα ΕΕ1--ΕΕ4 της \cref{sec:rq}, με τους αριθμούς του \cref{ch:evaluation}.}

\section{Περιορισμοί}
\label{sec:limitations}
\begin{itemize}
  \item Η ευθυγράμμιση εκτιμά στροφή μόνο γύρω από τον κατακόρυφο άξονα· προϋποθέτει σωστή κατακόρυφο από τη συσκευή.
  \item Η ακρίβεια περιορίζεται από το πόσο καλά ακουμπά ο χρήστης το σημείο με το χειριστήριο.
  \item Ο επανεντοπισμός των anchors εξαρτάται από τον φωτισμό και την υφή του χώρου.
  \item Ο ενσωματωμένος εξυπηρετητής προορίζεται μόνο για τοπικό δίκτυο (HTTP με κλειδί API).
  \item Η αναγνώριση φωνής εξαρτάται από υπηρεσία νέφους (Wit.ai).
  \item Η εφαρμογή χρησιμοποιεί το Built-in Render Pipeline της Unity και το SDK της Meta, άρα δεν μεταφέρεται αυτούσια σε άλλα headsets.
\end{itemize}

\section{Απειλές εγκυρότητας}
\label{sec:validity}
\guide{Μικρό δείγμα χρηστών, ένα εργοτάξιο, ένας τύπος συσκευής, ο ερευνητής είναι και ο δημιουργός του συστήματος (μεροληψία), εξοικείωση συμμετεχόντων με VR.}

\chapter{Συμπεράσματα και μελλοντική εργασία}
\label{ch:conclusions}

\guide{Έκταση 3--4 σελίδες. Καμία νέα πληροφορία· σύνοψη των συνεισφορών με τα κύρια αποτελέσματα.}

\section{Συμπεράσματα}
\label{sec:conclusions}
\todo{Σύνοψη.}

\section{Μελλοντική εργασία}
\label{sec:future}
\begin{itemize}
  \item \textbf{Αυτόματη επαναευθυγράμμιση} με ένα anchor ανά σημείο ελέγχου· η \cref{sec:exp-anchors} είναι η «Φάση~0» αυτής της πρότασης (\code{docs/proposals/auto-realignment.md}).
  \item \textbf{Διανύσματα μετακίνησης 3Δ και σύγκριση δύο ημερομηνιών}, με ένα endpoint «στιγμιότυπο έργου» στην πλατφόρμα (\code{docs/proposals/terraweb-data-for-new-features.md}, \code{displacement-visualization.md}).
  \item \textbf{Ευθυγράμμιση χωρίς σημεία ελέγχου}, με ταίριασμα του πλέγματος του χώρου που καταγράφει η συσκευή με το μοντέλο φωτογραμμετρίας (ICP \cite{besl1992icp}).
  \item \textbf{Συνεργατική χρήση}: κοινά anchors για πολλούς χρήστες στο ίδιο εργοτάξιο.
  \item \textbf{Φορητότητα}: OpenXR και Universal Render Pipeline για άλλα headsets.
\end{itemize}


% ---------------- Βιβλιογραφία ---------------------------------------
\printbibliography[heading=bibintoc,title={Βιβλιογραφία}]

% ---------------- Παραρτήματα ----------------------------------------
\appendix
% Τα παραρτήματα αριθμούνται Α, Β, … (ελληνικά) στο κείμενο· οι εσωτερικοί σύνδεσμοι του PDF θέλουν λατινικούς.
\renewcommand{\theHchapter}{app.\arabic{chapter}}% \Alph δίνει κι αυτό ελληνικά με το polyglossia
\chapter{Αναφορά τοπικού API}
\label{app:api}

Ο ενσωματωμένος εξυπηρετητής ακούει στη θύρα 8080 της συσκευής. Κάθε αίτημα, εκτός από το \code{/api/health}, απαιτεί το κλειδί της συσκευής στην επικεφαλίδα \code{X-Api-Key} ή \code{Authorization: Bearer}. Σώματα και απαντήσεις είναι JSON, εκτός αν σημειώνεται διαφορετικά.

\section{Endpoints}

\begin{longtable}{@{}p{4.6cm}p{1.6cm}p{8cm}@{}}
  \caption{Endpoints του τοπικού API.}\label{tab:api-full}\\
  \toprule
  Διαδρομή & Μέθοδος & Λειτουργία \\
  \midrule
  \endfirsthead
  \toprule
  Διαδρομή & Μέθοδος & Λειτουργία \\
  \midrule
  \endhead
  \code{/api/health} & GET & Κατάσταση και χρόνος λειτουργίας (χωρίς κλειδί) \\
  \code{/api/status} & GET & Σημεία, ευθυγράμμιση, συνεδρία \\
  \code{/api/logs} & GET & Πρόσφατα αιτήματα (κυκλικός buffer) \\
  \code{/api/coordinates} & GET, POST & Λίστα / προσθήκη σημείων \\
  \code{/api/coordinates/{id}} & DELETE & Διαγραφή ενός σημείου \\
  \code{/api/clear} & POST & Διαγραφή όλων των σημείων \\
  \code{/api/undo} & POST & Αναίρεση της τελευταίας αλλαγής σημείων \\
  \code{/api/layer} & POST & Φίλτρο ορατού στρώματος \\
  \code{/api/project} & POST & Φόρτωση έργου ευθυγράμμισης \\
  \code{/api/project/reset} & POST & Επαναφορά της ευθυγράμμισης του έργου \\
  \code{/api/calibrate/start} & POST & Έναρξη βαθμονόμησης με σημεία \\
  \code{/api/calibrate} & POST & Αποστολή πραγματικών συντεταγμένων σημείου \\
  \code{/api/calibrate/status} & GET & Κατάσταση βαθμονόμησης και υπόλοιπο RMS \\
  \code{/api/calibrate/reset} & POST & Επαναφορά βαθμονόμησης \\
  \code{/api/gps-origin} & POST & Ορισμός γεωγραφικής αρχής \\
  \code{/api/anchors} & GET & Κατάσταση του anchor αναφοράς \\
  \code{/api/session} & GET, POST & Εξαγωγή / εισαγωγή συνεδρίας \\
  \code{/api/session/new} & POST & Νέα ονομασμένη συνεδρία \\
  \code{/api/model} & POST & Μεταφόρτωση μοντέλου (OBJ, DXF ή zip) \\
  \code{/api/models} & GET & Λίστα μοντέλων στη συσκευή \\
  \code{/api/fieldnotes} & GET & Σημειώσεις πεδίου \\
  \code{/api/fieldnotes/photo/{file}} & GET & Φωτογραφία σημείωσης πεδίου (JPEG) \\
  \code{/api/reports}, \code{/api/reports/{file}} & GET & Αναφορές διαδρομών επιθεώρησης \\
  \code{/api/fieldtest}, \code{/api/fieldtest/{file}} & GET & Αρχεία CSV της δοκιμής anchors \\
  \bottomrule
\end{longtable}

\review{Ο πίνακας είναι από το README και τη λίστα διαδρομών του \code{LocalAPIServer.cs}. Έλεγξε τις μεθόδους του \code{/api/model} (μεταφόρτωση, διαγραφή) και του \code{/api/project} στον κώδικα, και αν υπάρχουν endpoints που λείπουν.}

\section{Κωδικοί απάντησης}

\begin{table}[htbp]
  \centering\small
  \caption{Κωδικοί σφάλματος.}
  \label{tab:api-errors}
  \begin{tabularx}{\textwidth}{@{}lX@{}}
    \toprule
    Κωδικός & Πότε \\
    \midrule
    400 & Κενό ή μη έγκυρο JSON, μη έγκυρες τιμές (NaN, εκτός ορίων) \\
    401 & Λείπει ή είναι λάθος το κλειδί \\
    405 & Γνωστή διαδρομή με λάθος μέθοδο \\
    408 & Το σώμα δεν διαβάστηκε εγκαίρως \\
    413 & Σώμα JSON πάνω από 512\,KB ή μοντέλο πάνω από 300\,MB \\
    429 & Υπέρβαση ορίου ρυθμού (60/λεπτό χωρίς κλειδί, 600/λεπτό με κλειδί) ή αποκλεισμός μετά από 10 λάθος κλειδιά· επικεφαλίδα \code{Retry-After} \\
    503 & Η εφαρμογή δεν απάντησε εγκαίρως (π.χ. συσκευή σε αναμονή) ή η ουρά εντολών είναι γεμάτη· επικεφαλίδα \code{Retry-After} \\
    \bottomrule
  \end{tabularx}
\end{table}

\section{Παράδειγμα: αποστολή σημείων}

\begin{lstlisting}[language=bash,caption={Αποστολή σημείων με cURL.}]
curl -X POST http://<quest-ip>:8080/api/coordinates \
     -H "Content-Type: application/json" \
     -H "X-Api-Key: <key>" \
     -d '{"coordinates":[{"id":"K5","x":12.4,"y":1.2,"z":-3.8,
          "label":"K5","color":"#FF0000","scale":0.1,
          "markerType":"pin","layer":"hazards"}],
          "clearExisting":false,"removeMissing":false}'
\end{lstlisting}

Τα πεδία \code{x}, \code{y}, \code{z} είναι σε μέτρα, με το \code{y} προς τα πάνω. Το \code{clearExisting} διαγράφει όλα τα σημεία πριν από την εισαγωγή, και το \code{removeMissing} διαγράφει όσα λείπουν από το αίτημα (συγχρονισμός). Αντί για \code{x/y/z}, ένα σημείο μπορεί να δοθεί με \code{latitude}, \code{longitude}, \code{altitude} και \code{"crs":"gps"}, αφού οριστεί πρώτα αρχή με \code{/api/gps-origin}.

\chapter{Πρωτόκολλο μελέτης χρηστών}
\label{app:userstudy}

\section{Έντυπο συναίνεσης}
\todo{Από την επιτροπή ηθικής του ιδρύματος, αν απαιτείται.}

\section{Ερωτηματολόγιο SUS}

Κλίμακα 1 (διαφωνώ απόλυτα) έως 5 (συμφωνώ απόλυτα) \cite{brooke1996sus}.

\guide{Χρησιμοποίησε δημοσιευμένη, σταθμισμένη ελληνική μετάφραση του SUS αν υπάρχει, και αναφέρ' την. Αλλιώς η παρακάτω είναι δική σου μετάφραση και πρέπει να δηλωθεί ως τέτοια.}

\begin{enumerate}
  \item Νομίζω ότι θα ήθελα να χρησιμοποιώ αυτό το σύστημα συχνά.
  \item Βρήκα το σύστημα άσκοπα πολύπλοκο.
  \item Νομίζω ότι το σύστημα ήταν εύκολο στη χρήση.
  \item Νομίζω ότι θα χρειαζόμουν τη βοήθεια τεχνικού για να μπορώ να χρησιμοποιήσω αυτό το σύστημα.
  \item Βρήκα ότι οι διάφορες λειτουργίες του συστήματος ήταν καλά ενσωματωμένες.
  \item Νομίζω ότι υπήρχε υπερβολική ασυνέπεια στο σύστημα.
  \item Φαντάζομαι ότι οι περισσότεροι άνθρωποι θα μάθαιναν να χρησιμοποιούν αυτό το σύστημα πολύ γρήγορα.
  \item Βρήκα το σύστημα πολύ δύσχρηστο.
  \item Ένιωσα πολύ σίγουρος/η χρησιμοποιώντας το σύστημα.
  \item Χρειάστηκε να μάθω πολλά πράγματα πριν μπορέσω να ξεκινήσω με αυτό το σύστημα.
\end{enumerate}

Βαθμολογία: για τα περιττά θέματα $(x-1)$, για τα άρτια $(5-x)$· το άθροισμα επί $2{,}5$ δίνει τιμή $0$--$100$.

\section{NASA-TLX}
\guide{Έξι κλίμακες 0--100: νοητικός φόρτος, σωματικός φόρτος, χρονική πίεση, απόδοση, προσπάθεια, απογοήτευση \cite{hart1988nasatlx}. Χρησιμοποίησε την απλή μορφή (raw TLX, μέσος όρος) εκτός αν ο επιβλέπων ζητήσει τις σταθμίσεις.}

\section{Εργασίες}
\begin{enumerate}
  \item Βρες τους αισθητήρες \fillin{Κ3}, \fillin{Κ7}, \fillin{Κ12} και ακούμπησέ τους.
  \item Ποιοι αισθητήρες είναι σε κατάσταση συναγερμού και πού βρίσκονται;
  \item Ποια περιοχή είχε τη μεγαλύτερη μετακίνηση τον τελευταίο μήνα;
\end{enumerate}

\section{Ημιδομημένη συνέντευξη}
\todo{5--6 ανοιχτές ερωτήσεις: τι βοήθησε, τι δυσκόλεψε, θα το χρησιμοποιούσες στη δουλειά, τι λείπει.}

\chapter{Ακατέργαστα αποτελέσματα}
\label{app:raw}

\section{Αρχεία δοκιμής πεδίου anchors}

Κάθε αρχείο CSV αντιστοιχεί σε ένα έργο και μία ημέρα, με μία γραμμή ανά anchor ανά μέτρηση. Οι στήλες είναι:

\begin{table}[htbp]
  \centering\small
  \caption{Στήλες των αρχείων \code{fieldtest/*.csv}.}
  \label{tab:fieldtest-columns}
  \begin{tabularx}{\textwidth}{@{}lX@{}}
    \toprule
    Στήλη & Σημασία \\
    \midrule
    \code{timestamp}, \code{session_id} & Χρονική στιγμή, εκκίνηση της εφαρμογής \\
    \code{project}, \code{point_id} & Έργο και σημείο ελέγχου \\
    \code{condition}, \code{site_changed} & Συνθήκη (sunny, cloudy, indoor, morning, noon, evening), αλλαγές εργοταξίου (yes/no) \\
    \code{app_version} & Έκδοση εφαρμογής \\
    \code{phase} & load, manual, touch, solution, create, undo, erase \\
    \code{anchor_uuid}, \code{anchor_age_h} & Anchor και ηλικία του σε ώρες \\
    \code{localized}, \code{seconds}, \code{tracked} & Επανεντοπίστηκε, σε πόσα δευτερόλεπτα, παρακολουθείται \\
    \code{dist_from_original_m} & Μετακίνηση σε σχέση με την αρχική θέση \\
    \code{errH_cm_*}, \code{errV_cm_*} & Οριζόντια/κατακόρυφη απόκλιση από τη λύση ή το νέο άγγιγμα \\
    \code{pair_errH_cm}, \code{pair_errV_cm} & Χειρότερη διαφορά απόστασης από τα άλλα anchors \\
    \code{headset_dist_m}, \code{note} & Απόσταση χρήστη, σημείωση \\
    \bottomrule
  \end{tabularx}
\end{table}

\todo{Συνοπτικοί πίνακες όλων των μετρήσεων. Τα αρχεία στο \code{thesis/data/}.}

\chapter{Εγκατάσταση και αναπαραγωγή}
\label{app:build}

\section{Εφαρμογή συσκευής}

\begin{enumerate}
  \item Unity 6 (6000.5) με υποστήριξη Android. Άνοιγμα του φακέλου \code{QuestMarkerApp/} από το Unity Hub.
  \item Το Meta XR Core SDK και τα υπόλοιπα πακέτα κατεβαίνουν αυτόματα από το \code{Packages/manifest.json}.
  \item Από το μενού \code{Quest Marker App} του Editor εφαρμόζονται οι ρυθμίσεις build για Quest (ελάχιστο API 32, στόχος API 34, IL2CPP, ARM64).
  \item Σύνδεση του Quest~3 με USB σε λειτουργία προγραμματιστή και \code{Build and Run}.
  \item Προαιρετικά, για φωνητικές εντολές: πακέτο \code{com.meta.xr.sdk.voice}, ρύθμιση Wit.ai και σύμβολο \code{META_VOICE_SDK}.
\end{enumerate}

\section{Εργαλεία υπολογιστή}

\begin{lstlisting}[language=bash,caption={Εγκατάσταση και έλεγχος των εργαλείων Python.}]
cd QuestMarkerApp/Server
python -m venv .venv
.venv\Scripts\activate          # Windows
pip install -r requirements.txt
pip install pymeshlab pillow pytest   # προετοιμασία μοντέλων, tests
python -m pytest -q                   # 127 tests
\end{lstlisting}

\section{Αναπαραγωγή των πειραμάτων}

\begin{enumerate}
  \item \code{git checkout thesis-v2} για την έκδοση της εφαρμογής που χρησιμοποιήθηκε στα πειράματα.
  \item Τα ακατέργαστα δεδομένα βρίσκονται στο \code{thesis/data/}.
  \item Τα σενάρια ανάλυσης στο \code{thesis/analysis/} (\code{pip install -r requirements.txt}) γράφουν πίνακες και γραφήματα στο \code{thesis/figures/generated/}.
  \item Η αναφορά μεταγλωττίζεται με \code{latexmk} (XeLaTeX + Biber) στον φάκελο \code{thesis/report/}.
\end{enumerate}


\end{document}
